% =============================================================================
%  《Chern--Simons 理论与 Jones 多项式》讲义 —— 导言区
%  编译：xelatex 两遍；中文字体 fandol
%  风格：连续叙事，公式嵌在正文，引用只进脚注
% =============================================================================
\documentclass[11pt,a4paper]{article}

\usepackage[UTF8, fontset=fandol, heading=true]{ctex}
\usepackage{amsmath,amssymb,mathtools,bm}
\usepackage[a4paper,left=2.7cm,right=2.7cm,top=2.8cm,bottom=2.8cm]{geometry}
\usepackage{graphicx}
\usepackage{xcolor}
\usepackage{enumitem}
\usepackage{booktabs}
\usepackage{array}
\usepackage[font=small,labelfont=bf,labelsep=quad]{caption}
\usepackage[colorlinks=true,
            linkcolor=blue!40!black,
            citecolor=blue!40!black,
            urlcolor=blue!40!black]{hyperref}

\linespread{1.18}
\setlength{\emergencystretch}{3em}
\allowdisplaybreaks[3]

\renewcommand{\figurename}{图}
\renewcommand{\tablename}{表}
\renewcommand{\refname}{参考文献}
\renewcommand{\contentsname}{目录}

\hypersetup{
  pdftitle={Chern-Simons 理论与 Jones 多项式},
  pdfauthor={WIT-1989-JONES 研讨班},
  pdfsubject={Chern-Simons theory, Jones polynomial, geometric quantization},
  pdfkeywords={Chern-Simons, Jones polynomial, conformal blocks, Verlinde formula},
  bookmarksnumbered=true
}

% ---- 数学符号 ----
\newcommand{\ii}{\mathrm{i}}
\DeclareMathOperator{\Tr}{Tr}
\DeclareMathOperator{\End}{End}
\DeclareMathOperator{\Hom}{Hom}
\DeclareMathOperator{\Inv}{Inv}
\DeclareMathOperator{\Vol}{Vol}
\newcommand{\dd}{\mathrm{d}}
\newcommand{\CC}{\mathbb{C}}
\newcommand{\ZZ}{\mathbb{Z}}
\newcommand{\RR}{\mathbb{R}}
\newcommand{\PP}{\mathbb{P}}
\newcommand{\cM}{\mathcal{M}}
\newcommand{\cL}{\mathcal{L}}
\newcommand{\cH}{\mathcal{H}}
\newcommand{\cO}{\mathcal{O}}
\newcommand{\cD}{\mathcal{D}}
\newcommand{\cN}{\mathcal{N}}
\newcommand{\cF}{\mathcal{F}}
\newcommand{\cR}{\mathcal{R}}

% ---- 脚注式引用：见原文某页 ----
\newcommand{\Wp}[1]{\textsuperscript{\,[W, p.#1]}}


\title{\bfseries Chern--Simons 理论、共形场论与 Jones 多项式\\[6pt]
\large 一份从三维量子化出发的入门讲义}
\author{WIT-1989-JONES 研讨班}
\date{2026 年}

\begin{document}

\maketitle

\begin{abstract}
\noindent
本讲义从一个具体问题出发：Jones 多项式是二维投影图上的组合量，为什么它能用三维量子场论来定义？
沿着这个问题，我们依次讲清楚 Chern--Simons 作用量、相空间与辛结构、几何量子化的三个步骤、
Wilson 线与穿刺点、Wess--Zumino--Witten 模型与共形块、手术计算与 Verlinde 公式，
最后落在两个常被混淆的对偶和纽结同调的范畴化上。

阅读本书只需要量子力学与一门场论或微分几何的基础；共形场论与几何量子化的必要内容都在正文里现讲。
我们的目标是让读者合上书之后，能自己把“三维理论量子化 $\to$ 模空间上的全纯截面 $\to$ 共形块 $\to$ 融合规则”
这条链完整复述一遍，并知道每一步在什么条件下成立。
\end{abstract}

\tableofcontents
\newpage

\section*{前言}
\addcontentsline{toc}{section}{前言}

\subsection*{这本书想回答什么}

1989 年，Witten 在《Communications in Mathematical Physics》上发表了一篇四十页的论文，
标题是《Quantum Field Theory and the Jones Polynomial》。论文的中心论断只有一句话：
Jones 多项式是一个三维量子场论的配分函数。这句话之所以重要，不是因为它给了多项式一个新的算法，
而是因为它把“纽结”从一个二维组合对象变成了三维拓扑的一部分，并由此把纽结理论、量子场论与
共形场论放进同一个框架。

本讲义试图把这句话讲成一个连贯的故事，而不是逐条注释一篇论文。我们只假定读者学过量子力学
和一学期的场论或微分几何；共形场论、几何量子化、仿射李代数这些工具，都会在用到的时候现讲。
读完本书，读者应当能独立地把这条链复述清楚：
\[
\begin{aligned}
&\text{三维 Chern--Simons 理论}
\;\longrightarrow\;
\text{曲面上的量子 Hilbert 空间}
\;\longrightarrow\;
\text{模空间上的全纯截面}\\
&\qquad\longrightarrow\;
\text{共形块}
\;\longrightarrow\;
\text{融合规则与 Verlinde 公式}.
\end{aligned}
\]

\subsection*{读者指南}

这本讲义按“一条主线、两个侧面”组织。主线是：Chern--Simons 理论在一个曲面上量子化，
得到的态空间恰好是 WZW 模型的共形块。两个侧面是：三维侧的 Wilson 线与手术计算，
二维侧的融合规则与 Verlinde 公式。读完八讲，读者应当能自己把这条主线复述一遍，
并知道每一步在什么条件下成立。

前置知识只有三样：量子力学（知道 Hilbert 空间与算符）、一门场论课（知道路径积分与规范场），
以及微分几何的基本语言（流形、丛、联络）。辛几何与复几何的概念在用到时现讲，
不需要提前准备。如果读者只想知道结论，读前言、第一讲与第八讲即可；
如果要动手算，第七讲是入口。

建议的读法是“两遍”：第一遍只读每讲开头的动机与结尾的小结，跳过所有带编号的推导；
第二遍再逐条推公式。书中标注 \Wp{n} 的页码是原文的印刷页，仅供核对，
不影响阅读顺序。

\subsection*{路线图}

全书分八讲。第一讲从 Jones 多项式的组合定义出发，说明它为什么需要一个三维的解释。
第二讲引进 Chern--Simons 作用量，并讲清楚它的两个反直觉性质：不依赖度规，以及耦合常数必须取整数。
第三讲把时空沿曲面切开，读出相空间：曲面上平坦联络的模空间。第四讲把相空间量子化，
这是全书技术上的核心，几何量子化的预量子化、极化、全纯截面三步都在这里。
第五讲把纽结放进来，讲 Wilson 线如何变成穿刺点上的一个约束，以及这个约束为什么不能经典地解。
第六讲转向二维侧，讲 WZW 模型、共形块与融合规则。
第七讲回到三维，用手术把配分函数算出来，并推出 Verlinde 公式。
第八讲讨论两个常见的误读——把本文当作全息对偶的雏形，或当作范畴化的预言——并给出仍然开放的问题。

\subsection*{约定}

\begin{itemize}[leftmargin=*]
  \item $M$ 表示紧致、定向、无边三维流形，$\Sigma$ 表示其中的曲面。$G$ 是紧致单 Lie 群，
        $d=\dim G$，$h^\vee$ 是它的对偶 Coxeter 数（$SU(N)$ 时 $h^\vee=N$）。
  \item 本文的 Chern--Simons 作用量取
        $L=\frac{k}{4\pi}\int_M\Tr(A\wedge\dd A+\frac23A\wedge A\wedge A)$，
        其中 $\Tr$ 取定义表示、长根平方为 $2$。$k$ 是整数。
  \item 引用原文时用边注 \Wp{370}，表示 Witten 论文的印刷页 370。
        教材页码在脚注里给出。
  \item 我们区分三种陈述：原文给出的构造、我们在正文里补出的推导、以及仍然只能算猜想的部分。
        正文中凡属后两类，都会明确写出“可验证的后果”，读者可以自己核对。
\end{itemize}

\section{从一个纽结多项式说起}

\subsection{Jones 多项式的定义}

设 $L\subset S^3$ 是一个定向链环，画成平面投影图。在每个交叉处，我们可以做三种局部操作：
把交叉换成两条平行线（0-光滑化）、换成另一种平行线（1-光滑化），或者保留交叉。
Jones 多项式由拆接关系（skein relation）刻画：
\begin{equation}
t^{-1}J(L_+)-t\,J(L_-)=(t^{1/2}-t^{-1/2})\,J(L_0), \tag{1.1}
\end{equation}
其中 $L_+,L_-,L_0$ 是三个投影图，除了一个交叉外完全相同，该交叉分别是正交叉、负交叉与光滑化。
再要求平凡纽结的值为 $1$，就唯一确定了链环的多项式 $J(L)\in\ZZ[t^{\pm1/2}]$。
例如三叶结的 Jones 多项式是
\[
J(\text{三叶结})=-t^{-4}+t^{-3}+t^{-1}.
\]

这个定义有几个漂亮的性质：它给出一个容易计算的算法；它把纽结的三种基本变换
（Reidemeister 变换）都吸收进了多项式；它还自然地带上一个变量 $t$，可以取任意复数。
但它也有一个根本的缺点：\textbf{定义依赖投影图}。为了证明 $J$ 是链环的不变量，
必须验证它在 Reidemeister 变换下不变；而在计算过程中，中间多项式来自完全不同的投影，
彼此之间没有任何内蕴的联系。

这正是 1980 年代纽结理论的核心困惑。人们知道 Jones 多项式存在、知道它很好用，
但不知道它“是什么”：它是不是某个更基本的、定义在三维空间里的量的影子？

\subsection{一个类比：1930 年代的量子电动力学}

Witten 在论文开头用了一个类比\Wp{352}。1930 年代，物理学家已经知道量子电动力学是正确的：
它给出与实验吻合的能级与散射截面。但当时没有人知道如何把 Maxwell 方程与量子力学
写成一种“明显相对论性”的形式；直到量子场论（特别是路径积分与协变微扰论）建立之后，
这个表述才出现。

Jones 多项式的处境类似：它的组合定义是“二维的”，而它最终应该来自某个“三维的”结构。
用 Witten 的话说，问题是：二维共形场论里究竟是哪个东西是三维的？

\subsection{二维的来源太多}

1980 年代中后期，纽结多项式从至少三条互不相干的路线冒出来\Wp{352}：
\begin{enumerate}[leftmargin=*]
  \item 可解格点模型与统计力学。某些二维格点模型的配分函数在临界点附近给出纽结不变量；
        相关的组合对象是 Temperley--Lieb 代数与 Yang--Baxter 方程。
  \item 二维共形场论的单值性。在有理共形场论里，共形块的单值性给出辫群表示，
        而辫群的迹与纽结多项式有关。
  \item 量子群与 Hopf 代数。Drinfeld 与 Jimbo 的量子群给出 Yang--Baxter 方程的解，
        进而给出纽结不变量。
\end{enumerate}
这三条路线各自都能算出 $J$，但它们彼此之间看不到明确的联系。一个自然的愿望是：
\textbf{为这种多样性找到一个统一的主题}。

\subsection{不变量家族与量子群}

Jones 多项式不是孤立的现象。1980 年代中后期，人们发现了一整族“量子不变量”：
HOMFLY 多项式（两个变量，来自 $SU(N)$ 的 $N$ 趋于一般值）、Kauffman 多项式（来自 $SO(N)$）、
以及更一般的 Reshetikhin--Turaev 不变量（来自量子群与模张量范畴）。
这些不变量有共同的结构：一个 $R$ 矩阵（Yang--Baxter 方程的解），配上辫群的表示，
再取迹。

这个观察很重要，因为它提示：纽结不变量不是某个二维组合技巧的产物，
而是某种“表示论 + 迹”的普遍结构。Witten 的工作给这个结构一个三维解释：
$R$ 矩阵来自 Wilson 线的融合，迹来自三维理论在曲面上的态空间。于是“量子群”
不再是代数上的巧合，而是三维拓扑场论的二维投影。

理解这一点，就能理解为什么后面要花大篇幅讲仿射李代数、共形块与融合规则：
这些看起来是二维的工具，实际上是在描述三维理论的态空间。

\subsection{一个不依赖度规的量子场论}

如果要用三维量子场论来定义纽结不变量，这个理论必须有一个非常特殊的性质：
它的配分函数不能依赖 $S^3$ 的度规，否则把纽结放进去时，结果会随度规改变。
换句话说，我们需要的是一种“拓扑场论”：定义在任意定向三维流形上、只依赖拓扑的量子场论。

这一点在广义相对论里有先例：引力是“对度规积分”的理论，它的定义不预设背景度规。
但广义相对论不是唯一的可能性。事实上，只要拉格朗日量本身不含度规，
量子化之后的配分函数就很可能不依赖度规。Witten 的观察是：三维的 Chern--Simons 理论
正是这样的理论，而它的 Wilson 线恰好给出纽结不变量\Wp{351}。

于是故事的目标变得清楚了：
\begin{enumerate}[leftmargin=*]
  \item 写出 Chern--Simons 作用量（第二讲）；
  \item 证明它不依赖度规、且耦合常数必须取整数（第二讲）；
  \item 把它量子化，得到曲面上的 Hilbert 空间（第三、四讲）；
  \item 把纽结写成 Wilson 线，得到一个三维不变量（第五讲）；
  \item 证明这个不变量就是 Jones 多项式，并算出来（第六、七讲）。
\end{enumerate}

这五步就是本讲义的主线。每一步都会遇到一个具体的技术问题：作用量如何变分、相空间是什么、
量子化用什么极化、穿刺点上的约束如何解、手术如何算。我们依次处理。

\subsection*{练习}
\begin{enumerate}[leftmargin=*]
  \item 用拆接关系 (1.1) 计算三叶结的 Jones 多项式，并与正文结果核对。
  \item 证明：若 $L$ 是平凡链环的 $n$ 个分量，则 $J(L)=(-t^{1/2}-t^{-1/2})^{n-1}$。
  \item 说明为什么 Jones 多项式不是链环的完全分类（给出两个不同纽结但同多项式的例子，或说明为什么难找）。
\end{enumerate}

\begin{quote}
\small
\textbf{本讲要点。} Jones 多项式有组合定义，但缺少内蕴的三维解释。
三维拓扑场论的配分函数（不依赖度规）与 Wilson 线的期望值，是可能的三维来源。
\end{quote}

\section{三维的规范场：Chern--Simons 作用量}

\subsection{联络与曲率}

设 $P\to M$ 是以紧致单 Lie 群 $G$ 为结构群的主丛，$A$ 是 $P$ 上的联络。在局部平凡化下，
$A=A_i\dd x^i$ 取值为 Lie 代数 $\mathfrak g$ 中。曲率是
\[
F=\dd A+A\wedge A,\qquad F_{ij}=\partial_iA_j-\partial_jA_i+[A_i,A_j].
\]
在规范变换 $A\mapsto A^{\,g}=g^{-1}Ag+g^{-1}\dd g$ 下，曲率协变地变换：$F\mapsto g^{-1}Fg$。
因此 $\Tr(F\wedge\star F)$ 这样的量是规范不变的。

我们还需要固定迹的归一化。取 $G$ 的定义表示，并要求
\[
\Tr(T^aT^b)=\tfrac12\delta^{ab},
\]
即长根平方为 $2$ 的归一化。这个选择决定了下面所有公式里 $k$ 的单位；换表示或换归一化，
$k$ 的整性条件会跟着换。讲义附录里有一张各教材约定对照表，读者在比较不同书时务必先对齐这一项。

\subsection{为什么不用 Yang--Mills 作用量}

最自然的规范场作用量是 Yang--Mills 作用量
\[
L_{\mathrm{YM}}=\int_M\sqrt g\,g^{ik}g^{jl}\Tr(F_{ij}F_{kl})\,\dd^3x.
\]
它需要度规：$\sqrt g$ 是体积元，$g^{ik}g^{jl}$ 把两个曲率指标的上下位置缩并。
这两处对度规的依赖性质不同：前者是体积元，换度规只是改标度，可以用“不积分体积”的方式绕开；
后者把不同分量混合起来，\textbf{绕不开}——换一个度规，作用量里的混合方式就变了。

既然我们要的是一个不依赖度规的配分函数，就不能用 Yang--Mills。需要另一个只依赖微分拓扑的作用量。

\subsection{Chern--Simons 作用量}

在三维流形上，存在一个只用到 $A$ 与 $\dd$ 的三次型：
\begin{equation}
L=\frac{k}{4\pi}\int_M\Tr\left(A\wedge\dd A+\frac23A\wedge A\wedge A\right)
=\frac{k}{8\pi}\int\dd^3x\,\varepsilon^{ijk}\Tr\left(A_i\partial_jA_k+\frac23A_i[A_j,A_k]\right). \tag{2.1}
\end{equation}
这就是 Chern--Simons 作用量。它有三个决定性的性质。

\textbf{第一，它只依赖定向，不依赖度规。} 被积式是三维流形上的 $3$-形式，积分只需要定向，
不需要 $\sqrt g$，也不需要缩并。

\textbf{第二，它的变分给出曲率。} 直接计算给出
\begin{equation}
\delta L=\frac{k}{2\pi}\int_M\Tr(\delta A\wedge F). \tag{2.2}
\end{equation}
因此经典运动方程是 $F=0$：联络必须是平坦的。这个结果的物理含义值得停下来想一想：
三维 Chern--Simons 理论没有传播的自由度，它的经典解是“平坦联络”，
而不是像 Yang--Mills 那样有波。这正是它能成为拓扑场论的原因，也是第三讲相空间为有限维的根源。

\textbf{第三，它在小规范变换下不变，在大规范变换下平移 $2\pi$ 的整数倍。} 对小规范变换
$g$ 接近恒等，$\delta L=0$；对不平凡的大规范变换，$L$ 改变一个常数。这个常数决定了 $k$ 必须取整数。

下面把这三条依次算清楚。

\subsection{变分}

把 (2.1) 代入 (2.2) 并分部积分，可以把 $F$ 提出来。规范场论里常用的恒等式是
\[
\delta\int\Tr(A\wedge\dd A)=\int\Tr(\delta A\wedge\dd A)-\int\Tr(A\wedge\dd\,\delta A)
=2\int\Tr(\delta A\wedge\dd A),
\]
最后一步用了 $\dd A=-A\wedge\dd$（在三维、对一次形式的 $A$）与迹的循环性；
而 $\delta\frac23\int\Tr(A^3)=2\int\Tr(\delta A\wedge A\wedge A)$。
两式相加即得 (2.2)。

\subsection{大规范变换与 $k$ 的整性}

设 $g:M\to G$ 是规范变换，其同伦类由绕数 $m\in\ZZ$ 刻画（因为 $\pi_3(G)=\ZZ$ 对紧致单群成立）。
在 $A\mapsto g^{-1}Ag+g^{-1}\dd g$ 下，经过一段计算，$L$ 的变化为
\begin{equation}
L\longmapsto L+\frac{k}{4\pi}\cdot\frac13\int_M\Tr\left((g^{-1}\dd g)^3\right)
=L+\frac{k}{4\pi}\cdot\frac13\cdot 24\pi^2m
=L+2\pi k\,m, \tag{2.3}
\end{equation}
其中用到了恒等式
\[
m=\frac{1}{24\pi^2}\int_M\Tr\left((g^{-1}\dd g)^3\right)\in\ZZ.
\]
在量子理论里，物理量必须是 $\exp(\ii L)$ 的期望值，而这个量在规范变换下必须不变。
因此要求 $2\pi k m$ 是 $2\pi$ 的整数倍，即
\begin{equation}
k\in\ZZ. \tag{2.4}
\end{equation}
这是 Chern--Simons 理论最著名的性质：耦合常数被量子化。它的根源是 $\pi_3(G)$ 的非平凡性，
与二维共形场论里仿射李代数的中心荷取整数是同一件事的两面。

\subsection{Wilson 线与配分函数}

纽结进入理论的方式是 Wilson 线。给定链环 $C$ 与它的一个表示 $R$，定义
\begin{equation}
W_R(C)=\Tr_R\,P\exp\int_C A_i\dd x^i, \tag{2.5}
\end{equation}
其中 $P$ 表示路径有序。这个量不需要度规，只用到 $A$ 沿曲线的积分，因此是拓扑场论里自然的可观测量。
更一般地，对链环 $C_i$ 与表示 $R_i$，配分函数定义为
\begin{equation}
Z(M;C_i,R_i)=\int\cD A\,\exp(\ii L)\prod_i W_{R_i}(C_i). \tag{2.6}
\end{equation}
这里的路径积分不是严格定义的，但它的半经典展开和手术计算都可以严格化——这是第七讲的内容。

Witten 的中心论断是：当 $G=SU(N)$、所有 $R_i$ 取定义表示时，$S^3$ 上的这个配分函数
恰好给出 Jones 多项式，而两个参数 $N$ 与 $k$ 合起来给出二变量推广\Wp{356}。
当 $G=SO(N)$、$R$ 取 $N$ 维表示时，得到的是 Kauffman 多项式。

\subsection{弱耦合是大 $k$}

最后提醒一个容易混淆的地方。Chern--Simons 理论的“弱耦合极限”是 $k\to\infty$，
不是 $k\to0$。原因是 $k$ 在作用量里是一个整体因子：$k$ 越大，路径积分的相位振荡越剧烈，
半经典近似越好。这与 Yang--Mills 的 $1/g^2$ 展开正好相反，读者在比较不同文献时要特别小心。

\subsection{一个几何解释：Chern--Simons 是四维示性类的边界项}

$k$ 为什么必须是整数，还可以从一个更几何的角度看。在四维流形 $X$ 上，
$p_1=\frac{1}{8\pi^2}\Tr(F\wedge F)$ 是一个整上同调类（由 Atiyah--Singer 指标定理）。
它的一个陈--西蒙斯形式是
\[
\dd\,\omega_{\mathrm{CS}}(A)=\Tr(F\wedge F),
\]
而 (2.1) 的被积式就是 $\omega_{\mathrm{CS}}$ 的积分。如果 $X$ 以 $M$ 为边界，
那么 $\int_X\Tr(F\wedge F)=\int_M\omega_{\mathrm{CS}}$。左边是整数（超荷/瞬子数），
所以右边在 $M$ 上必须整——这正是 $k\in\ZZ$ 的几何来源。
换句话说，Chern--Simons 理论是四维拓扑不变量（第二陈类、瞬子数）在三维的“边界理论”。

这个观点也解释了为什么三维是 Chern--Simons 的“自然维度”：三次型需要一个三维积分，
而它的微分是四维示性类。更高维的推广是 $2n-1$ 维的 Chern--Simons 理论，
但物理上最相关的是 $n=2$。

\subsection{归一化、level 与 $q$}

同一篇论文里的 $k$，在不同的书里可能差一个因子。原因是 Lie 代数的迹有多种归一化：
有人取定义表示、长根平方为 $2$；有人取 Killing 型；有人用“中心荷归一到 $1$”的约定。
本节统一用定义表示、$\Tr(T^aT^b)=\frac12\delta^{ab}$；在这个约定下 (2.1) 的 $k$ 与仿射代数的 level 一致。

量子化之后常用的变量是
\[
q=\exp\left(\frac{2\pi\ii}{k+h^\vee}\right),
\]
其中 $h^\vee$ 是对偶 Coxeter 数（$SU(N)$ 时为 $N$）。$q$ 的单位根性质是“可积表示只有有限多个”
的直接体现：当 $q$ 是单位根时，量子群的表示范畴退化，只剩下有限多个可积对象。
Jones 多项式里的 $t$ 与这个 $q$ 的关系是 $t=q^{2}$（不同文献的约定略有差别），
读者在比较时必须先对齐 $k$、$h^\vee$ 与 $q$ 的定义。

\subsection{路径积分可算吗}

(2.6) 的路径积分不是严格定义的对象：三维流形上的规范场测度没有已知的构造。
但 Chern--Simons 理论有一个特殊性质：它的路径积分只依赖拓扑，
因此在半经典极限（$k\to\infty$）与手术分解下都可以“算”。Witten 的做法是：
把三维流形沿曲面切开，用曲面上的态空间做粘合；曲面上的态空间是有限维的，
所以粘合是一个有限维的线性代数操作。这样，路径积分就化归为一组自洽的有限维数据
（$S$ 矩阵、融合系数、标架相位），而这些数据由二维共形场论提供。

换句话说，这个理论不是靠“把路径积分算出来”来定义的，而是靠一组公理（TQFT 公理）
加上一组由共形场论给出的解来定义的。第七讲会把这组数据写具体。

\subsection{带边界的情形与边界项}

前面都假设 $M$ 无边。一旦 $M$ 有边界 $\partial M=\Sigma$，变分 (2.2) 会多出一个边界项：
\begin{equation}
\delta L=\frac{k}{2\pi}\int_M\Tr(\delta A\wedge F)
-\frac{k}{4\pi}\int_\Sigma\Tr(\delta A\wedge A). \tag{2.7}
\end{equation}
要让变分良定义，必须指定边界上 $A$ 的什么分量固定。常见的选择是固定 $A$ 在 $\Sigma$ 上的
一个复结构分量（即固定 $A^{(0,1)}$）；这时边界项消失，体内方程仍是 $F=0$。

这个边界项在故事里很重要。第五讲把纽结写成穿刺点时，穿刺点附近的小圆盘就是带边界的曲面；
边界项决定了“在边界上插入什么态”。第七讲的手术公式，本质上也是把边界上的态求和掉。
换句话说，Chern--Simons 理论不是“只活在三维内部”的理论：它的边界携带一个二维的态空间，
这正是它能与共形场论对接的技术原因。

\subsection{$U(1)$ 情形：一个可以完全算清的例子}

把群换成 $U(1)$ 时，一切都可以显式写出。$A$ 是一个普通一形式，$F=\dd A$，作用量是
\[
L=\frac{k}{4\pi}\int_M A\wedge\dd A .
\]
运动方程 $\dd A=0$，即 $A$ 是闭形式。把 $A$ 分解成“恰当部分 + 调和部分”，
恰当部分可以规范掉，只剩下调和部分；于是理论化归为有限维的辛向量空间上的高斯积分。
在曲面 $\Sigma$ 上，$U(1)$ 理论的 Hilbert 空间是 theta 函数空间，维数为 $k^g$；
环面上（$g=1$）维数为 $k$，与第四讲的截面计数一致。

$U(1)$ 情形里 $k$ 的整性同样来自大规范变换的绕数，与第二讲完全平行；
但这里没有非阿贝尔的复杂性，所以可以拿来检验每一个公式。读者若对某一步不放心，
先退到 $U(1)$ 算一遍，通常能看清结构。

\subsection{微扰展开：$1/k$ 展开给出什么}

$k\to\infty$ 时，路径积分可以用驻点法展开，展开参数是 $1/k$，展开的图形是三维费曼图。
有一个漂亮的结果：这些费曼图给出的不变量，恰好是 Vassiliev 不变量（有限型不变量）的各级；
换句话说，Chern--Simons 微扰论给出纽结不变量的一个无穷族，Jones 多项式只是其中最粗的一层。

这件事有两个意义。第一，$1/k$ 展开不是形式幂级数，而是有几何意义的：
每一项对应一类奇点的消解。第二，它说明为什么 Jones 多项式不够：
有限型不变量的层级比单个多项式丰富得多。

不过微扰展开有一个技术困难：三维费曼图会出现紫外发散，需要正则化；
而正则化会引入标架依赖，这正是第七讲标架反常的来源。也就是说，
微扰论的自洽性要求标架结构，而标架在拓扑场论里本来就是一个必须处理的数据。

\subsection{为什么三维 Chern--Simons 是可解的}

一个三维量子场论通常无法精确求解：自由度太多，相互作用太强。Chern--Simons 之所以例外，
原因可以一句话概括：它没有局域自由度。

作用量 (2.1) 不含度规，也不含 $F$ 的平方项，所以它不是一个“传播”的理论：
没有引力子、没有规范玻色子的传播子，唯一的动力学是整体（拓扑）的。
把流形沿曲面切开后，每一侧的态空间是有限维的；所有的“物理”都集中在边界上。
这就是为什么它可以用二维共形场论的数据完全描述。

这个观察在技术上体现为两点。第一，相空间是有限维的（第三讲）：它不是“场的空间”，
而是“解的空间”。第二，配分函数只依赖流形的拓扑（手术不变量，第七讲）。
物理上，这类理论叫拓扑场论（TQFT）；Witten 1988 年的另一篇文章给出了 TQFT 的公理，
而本文是它最成功的例子。

从另一个角度看，Chern--Simons 与二维模型的关系是“同一组数据的两种组织方式”：
三维侧是拓扑的、整体的，二维侧是共形的、局域的。这个对偶关系不是本文发明的，
但本文第一次说明它为什么成立：三维理论在曲面上的态空间，就是二维理论的共形块。

\subsection*{练习}
\begin{enumerate}[leftmargin=*]
  \item 验证 (2.2)：从 (2.1) 出发分部积分，注意三维中 $A\wedge\dd A$ 与 $A\wedge A\wedge A$ 的系数 $\frac23$ 如何分配。
  \item 证明 $m=\frac{1}{24\pi^2}\int\Tr((g^{-1}\dd g)^3)$ 对 $g:S^3\to SU(2)$ 取整数，并计算 $g$ 为恒等映射的倍数时的值。
  \item 对 $G=SU(2)$，用 $\Tr(T^aT^b)=\frac12\delta^{ab}$ 写出作用量在 $A=A^aT^a$ 下的显式形式，并检验它与 (2.1) 的系数一致。
  \item 说明为什么 $k\to0$ 不是弱耦合极限。
\end{enumerate}

\begin{quote}
\small
\textbf{本讲要点。} Chern--Simons 作用量只依赖定向；(2.2) 给出 $F=0$ 的经典解；
大规范变换 (2.3) 强制 $k\in\ZZ$；Wilson 线 (2.5) 与配分函数 (2.6) 是纽结不变量的候选；
$k$ 的几何来源是四维示性类的边界项。
\end{quote}

\section{沿曲面切开：相空间是平坦联络}

\subsection{把时空切成两半}

量子化要从正则形式出发，而正则形式需要一个“时间”方向。三维流形 $M$ 没有天然的
时间，但我们可以沿一个闭曲面 $\Sigma$ 把它切开，写成 $\Sigma\times\RR$。
在 $\Sigma\times\RR$ 上，联络分解为
\[
A=A_0\,\dd t+A_i\,\dd x^i,
\]
其中 $A_0$ 是沿时间方向的联络分量（它是 $\mathfrak g$ 值的函数，不是一形式）。

把 (2.1) 代入并显式写出 $3$-形式的各项，注意 $\dd t\wedge\dd x^i\wedge\dd x^j$ 中
$i,j$ 只取两个空间方向，三次项只剩下 $A_0[A_i,A_j]$ 的组合，而 $A_0\wedge A_i\wedge A_j$
的符号使它与 $A_i\partial_tA_j$ 合并。取 $A_0=0$ 的规范，作用量变成
\begin{equation}
L=\frac{k}{8\pi}\int\dd t\int_\Sigma \varepsilon^{ij}\Tr\left(A_i\,\partial_tA_j\right)\dd^2x. \tag{3.1}
\end{equation}
这是一个对时间一阶的作用量：没有 $\dot A^2$ 项。它看起来像
$\int p\,\dot q$，其中 $p$ 就是 $A$ 的某个分量。因此相空间就是“所有 $A_i$”构成的空间，
并且它的辛形式可以立刻读出来。

\subsection{泊松括号}

从 (3.1) 读出正则对易关系。把 $\Sigma$ 上的两个分量记作 $A_1,A_2$（即 $A_i$ 的两维指标），
对 Lie 代数指标 $a,b$，有
\begin{equation}
\{A_i^a(x),A_j^b(y)\}=\frac{4\pi}{k}\,\varepsilon_{ij}\,\delta^{ab}\,\delta^2(x-y). \tag{3.2}
\end{equation}
注意括号里的 $4\pi/k$：耦合常数出现在辛形式的分母里。这正是为什么
“$k\to\infty$ 是弱耦合”在相空间里看得最清楚——$k$ 越大，$\hbar_{\text{eff}}=4\pi/k$ 越小。

\subsection{高斯定律与辛商}

$A_0$ 没有出现在 (3.1) 里，它是一个拉格朗日乘子。对 $A_0$ 变分给出约束：
\begin{equation}
\varepsilon^{ij}F_{ij}^a=0. \tag{3.3}
\end{equation}
在二维曲面上，曲率 $F_{ij}$ 只有一个独立分量（$F_{12}=-F_{21}$），所以 (3.3) 就是
\[
F=0,
\]
即联络必须平坦。这与第二讲从三维变分得到的结论一致：Chern--Simons 理论的经典解是平坦联络。

(3.3) 是“第一类约束”：它生成规范变换，因此不能简单地令其为零，而要按规范群取商。
用几何力学的语言，这是一个\textbf{辛商}（Marsden--Weinstein 约化）：先限制到约束面 $F=0$，
再按规范群的作用取商。两步合起来给出相空间
\begin{equation}
\cM=\left\{\text{$\Sigma$ 上的平坦联络}\right\}\big/\left\{\text{规范变换}\right\}
\cong\Hom(\pi_1\Sigma,G)/G. \tag{3.4}
\end{equation}
最后这个等式说明：平坦联络由它的和乐表示（$\pi_1\Sigma\to G$ 的同态）决定，只差一个共轭。

\subsection{相空间的维数}

$\Sigma$ 的亏格为 $g\ge2$ 时，$\pi_1\Sigma$ 有 $2g$ 个生成元，$\Hom(\pi_1\Sigma,G)$ 的维数是 $2gd$。
除以 $G$ 的共轭作用（$d$ 维），得到
\begin{equation}
\dim_\RR\cM=(2g-2)d,\qquad g\ge2. \tag{3.5}
\end{equation}
这个公式对 $g\ge2$ 成立。$g=1$ 时情形不同：一般位置的平坦联络有更大的稳定化子
（它的和乐落在某个极大环面里），所以不能简单地减去 $d$。对 $SU(2)$，环面上的模空间是
“枕头套” $T^2/\{\pm1\}$，实维数为 $2$。$g=0$ 时平坦联络只有平凡解，$\cM$ 是一个点。

(3.5) 是本书的第一个定量结论：相空间是有限维的，而且维数只依赖亏格与群，
不依赖任何连续参数。半经典地，量子 Hilbert 空间的维数应当约为
$\mathrm{Vol}(\cM)/(2\pi\hbar)^{n}$（$2n=\dim\cM$）；因为 $\cM$ 紧致、辛体积有限，这个数有限。
于是\textbf{Chern--Simons 理论在曲面上的 Hilbert 空间是有限维的}——这在三维量子场论里非常罕见，
也是整个故事能算得动的原因。

\subsection{复结构：从平坦联络到全纯丛}

要量子化，还需要给 $\cM$ 一个复结构。$\Sigma$ 上选一个复结构 $J$（等价地，一个亏格 $g$
的黎曼面结构）。Narasimhan--Seshadri 定理说：$\Sigma$ 上的平坦 $G$-联络与
“稳定的全纯 $G_\CC$-丛”之间存在自然的对应——一个平坦联络可以分解为
“全纯结构 + 与它相容的平坦幺正结构”。在这个对应下，$\cM$ 成为复流形 $\cM_J$，
而且它带有 Kähler 结构（对 $SU(N)$，$\cM_J$ 是 $c_1=0$ 的稳定秩 $N$ 全纯丛的模空间），
还是一个射影代数簇。

这个 Kähler 结构是下一讲量子化的入口：有了复结构，就能谈“全纯截面”，
而全纯截面正是我们要的量子态。

\subsection{辛商的具体形式：Atiyah--Bott}

约束面 $F=0$ 上的辛形式可以从 (3.2) 直接读出。设 $\delta A$ 是 $A$ 的一个变分，
对应切向量；辛形式是
\begin{equation}
\omega(\delta_1A,\delta_2A)=\frac{k}{4\pi}\int_\Sigma\Tr\left(\delta_1A\wedge\delta_2A\right). \tag{3.6}
\end{equation}
它显然是闭的（因为它是常系数二次型），并且在下述意义下非退化：若 $\delta_1A$ 与所有
满足 $F=0$ 的 $\delta_2A$ 的内积为零，则 $\delta_1A$ 必须是“纯规范”的。

Atiyah--Bott 的观察是：这个辛形式可以等价地写成模空间上的一个 Kähler 形式。
在 $\cM_J$ 上，切空间可以等同于 $H^1(\Sigma;\mathfrak g)$（一阶形变），
而 $\omega$ 就是 $H^1$ 上的 cup 积配对。于是 $\cM_J$ 的辛结构、复结构、
以及 (3.6) 的 Kähler 形式三者相容。这个等价性是第四讲能谈“Kähler 极化”的前提。

从 (3.6) 还可以看出 $k$ 的作用：$\omega$ 正比于 $k$。因此 $k$ 越大，辛体积越大，
半经典极限越好，量子 Hilbert 空间的维数越大——这与第四讲的 Weyl 律一致。

\subsection{小结}

把三维理论沿曲面切开之后，我们得到了一个具体得多的对象：有限维、紧致、Kähler 的模空间
$\cM_J$。三维量子场论的 Hilbert 空间，就是这个模空间上某个线丛的全纯截面空间。
剩下的问题是：哪个线丛？为什么是全纯截面？这是第四讲的主题。

\subsection*{练习}
\begin{enumerate}[leftmargin=*]
  \item 由 (3.1) 直接读出 (3.2) 的泊松括号（提示：把作用量写成 $\int p\dot q$ 的形式）。
  \item 计算 $g=0$ 与 $g=1$ 情形下 $\cM$ 的维数，并解释 $g=1$ 时公式 (3.5) 为什么不适用。
  \item 证明：$\cM$ 上由 (3.6) 定义的辛形式在下述意义下非退化——若它与所有闭形式的内积为零，则该形式是恰当的。
\end{enumerate}

\begin{quote}
\small
\textbf{本讲要点。} 正则形式给出辛形式与泊松括号 (3.2)；高斯定律 (3.3) 经辛商给出
平坦联络模空间 (3.4)，维数 $(2g-2)d$；选复结构后 $\cM_J$ 是 Kähler 模空间。
\end{quote}

\section{量子化：预量子化、极化与全纯截面}

\subsection{为什么量子化不平凡}

我们现在有了一个有限维、紧致、Kähler 的相空间 $\cM_J$，要把它量子化。
如果 $\cM_J$ 是 $\RR^{2n}$，量子化就是“把 $p$ 换成 $-i\hbar\partial_x$”，没有歧义。
但 $\cM_J$ 是紧致的，没有整体的“位置与动量分离”：任何试图把坐标分成 $q$ 与 $p$ 的做法
都只能局部进行。更具体地，量子态应当是相空间上的“函数”，但“函数”是什么空间？
如果取所有 $L^2$ 函数，维数是无限的，而物理上我们希望有限维（因为相空间体积有限）。

几何量子化给出了一套系统的做法，分三步：\textbf{预量子化}、\textbf{极化}、\textbf{半形式修正}。
这三步分别解决“线丛从哪来”“态是什么”“能级怎么移”三个问题。

\subsection{第一步：预量子化线丛}

相空间 $(\cM,\omega)$ 的预量子化，是找一个复线丛 $L\to\cM$，带一个酉联络 $\nabla$，
使得它的曲率是
\[
F_\nabla=-\ii\,\omega
\]
（这里 $\omega$ 按“每单位体积一个态”的归一化，见下）。这样的线丛存在的充要条件是
\begin{equation}
\frac{[\omega]}{2\pi}\in H^2(\cM;\ZZ). \tag{4.1}
\end{equation}
这就是“整性条件”。满足它之后，$L$ 上的截面就构成预量子 Hilbert 空间。

在我们的情形里，$\cM$ 上的辛形式由 (3.2) 给出：
\[
\omega=\frac{k}{4\pi}\int_\Sigma\Tr\left(\delta A\wedge\delta A\right).
\]
注意 $\omega$ 本身正比于 $k$。整性条件 (4.1) 于是恰好要求 $k$ 取整数——这与第二讲从大规范变换
得到的 $k\in\ZZ$ 是同一件事的两种说法。这一点很重要：\textbf{量子化的自洽性自动包含了耦合常数的量子化}。

那么线丛 $L$ 具体是什么？答案是行列式线丛。$\cM$ 是平坦联络的模空间，而平坦联络可以看成一族
椭圆算子（$\bar\partial_A$）的参数的模空间；这族算子的行列式构成 $\cM$ 上的一个线丛，
Quillen 给了它一个自然的度量与联络，其曲率正是上面的 $\omega$。这个构造的好处是：
\textbf{线丛是先验存在的}，所以整性条件自动满足，不需要额外验证。

\subsection{第二步：极化}

预量子 Hilbert 空间是 $L$ 的 $L^2$ 截面空间。它太大了：$\cM$ 的维数是 $2n$ 时，
这个空间是“$n$ 个坐标的函数”，而量子态应该只依赖一半的变量。
物理上，这对应“不能同时指定位置与动量”。

解决办法是选一个\textbf{极化}：复化切丛 $T\cM\otimes\CC$ 的一个子丛 $P$，满足
\begin{enumerate}[leftmargin=*]
  \item $P$ 是拉格朗日的：$\omega|_P=0$；
  \item $P$ 是可积的：对合（即 $P$ 的截面在 Lie 括号下封闭）；
  \item $\dim_\CC P=n$。
\end{enumerate}
量子态定义为“沿 $P$ 协变常数的截面”：对 $P$ 的任意截面 $X$，要求 $\nabla_Xs=0$。
这样，态只依赖 $P$ 的叶空间上的变量，维数就从无限降下来。

最常见的两种极化是：
\begin{itemize}[leftmargin=*]
  \item \textbf{实极化}：$P$ 是实切丛的复化，对应“位置空间”或“动量空间”。在 $\RR^{2n}$ 上取
        $P=\mathrm{span}\{\partial/\partial p_i\}$，得到的态只依赖 $x$。
  \item \textbf{Kähler 极化}：$P=T^{0,1}\cM$。此时“沿 $P$ 协变常数”就是“全纯”，
        量子态是全纯截面。
\end{itemize}
我们的 $\cM_J$ 是 Kähler 的，自然取 Kähler 极化。于是
\begin{equation}
\cH_\Sigma=H^0(\cM_J,\,L^{\otimes k}), \tag{4.2}
\end{equation}
即线丛 $L$ 的 $k$ 次幂的全纯截面空间。这里出现 $L^{\otimes k}$ 而不是 $L$，
是因为 $\omega$ 正比于 $k$：整性条件要求线丛的曲率是 $\omega$，而 $\omega=k\omega_0$，
所以线丛取 $L^{\otimes k}$，其中 $L$ 的曲率是 $\omega_0$。

\subsection{维数与半经典极限}

(4.2) 的维数由 Riemann--Roch（或指标定理）给出。对 $SU(N)$、亏格 $g$ 的情形，
$k\to\infty$ 时
\[
\dim H^0(\cM_J,L^{\otimes k})\sim \frac{\mathrm{Vol}(\cM)}{(2\pi)^n n!}\,k^n,
\]
即“每单位体积一个态”的 Weyl 律，其中 $2n=\dim\cM$。

最简单地，取 $g=1$、$G=SU(2)$：此时 $\cM_J$ 是枕头套 $T^2/\{\pm1\}$，
而 $\dim H^0(\cM_J,L^{\otimes k})=k+1$——这正是环面上的 theta 函数空间维数。
$k=1$ 时给出 $2$，$k=2$ 时给出 $3$，与第八讲算 $S^3$ 配分函数时用到的数一致。

值得把 $g=1$ 的情形算清楚，因为它包含了几何量子化最朴素的所有要素。取环面
$\Sigma=\CC/(\ZZ+\tau\ZZ)$。度为 $k$ 的全纯线丛的标准截面由 theta 函数给出：
\[
\vartheta_{j}(z)=\sum_{m\in\ZZ}\exp\left(\pi\ii k(m+\tfrac jk)^2\tau
+2\pi\ii k(m+\tfrac jk)z\right),\qquad j=0,1,\dots,k-1,
\]
共 $k$ 个。但正确的量子化还要乘上半形式因子 $K^{1/2}$；对 $SU(2)$ 它是 $L^{-2}$，
于是实际空间是 $H^0(L^{k}\otimes K^{1/2})$，维数为 $k+1$。这正好等于
“有效 level $k+2$ 的可积表示个数减一”，与 (4.3) 的平移一致。
这个例子说明：半形式不是一个可选的修饰，它改变维数。

\subsection{第三步：与复结构无关吗？}

一个自然的问题是：$\cH_\Sigma=H^0(\cM_J,L^{\otimes k})$ 依赖复结构 $J$，
但物理上 Hilbert 空间不应该依赖“怎么切”。这两个说法如何统一？

答案是：当 $J$ 变化时，$H^0(\cM_J,L^{\otimes k})$ 拼成 Teichmüller 空间上的一个向量丛，
而这个丛带有\textbf{射影平坦}联络（Hitchin 联络）。所谓射影平坦，是指平行移动只在一个
整体相位上有歧义；沿闭路平行移动给出一个射影变换，其单值性正是映射类群在 $\cH_\Sigma$ 上的作用。
因此严格地说，$\cH_\Sigma$ 不是一个“与 $J$ 无关的向量空间”，而是一个“射影平坦丛的纤维”；
只有射影化之后的 $\PP\cH_\Sigma$ 才是典范的。

这一点在第五讲之后会变得关键：共形块正是这个射影平坦丛的纤维，
而映射类群的作用给出辫群表示，从而给出纽结不变量。

\subsection{半形式修正与能级平移}

还有第三步：半形式修正。Kähler 极化要求我们取 $L^{\otimes k}$ 的截面，
但严格地，量子化还需要一个“半形式”因子 $K^{1/2}$（$K$ 是 $\cM_J$ 的典范丛）。
对 $\cM_J$ 而言，$K$ 与 $L$ 的某个幂有关：$K=\mathrm{Det}^{-2h^\vee}$，
因此 $K^{1/2}=L^{-h^\vee}$。把这个因子乘进去，等价于把 $k$ 换成
\begin{equation}
k\;\longrightarrow\;k+h^\vee. \tag{4.3}
\end{equation}
对 $SU(N)$，$h^\vee=N$，所以 $k\to k+N$；对 $SU(2)$，$k\to k+2$。

(4.3) 是一个“能级平移”。它在物理上来自标架反常，在数学上来自半形式。
Witten 的论文在正文里只用 $L^{\otimes k}$，把这个修正略去了；但所有需要
$q=\exp(2\pi\ii/(k+N))$ 的公式，都已经隐含了它。读者在比较不同文献时，
必须注意对方的 $k$ 是“裸”的还是“已平移”的。

\subsection{Hitchin 联络}

“射影平坦”这句话可以写成一个具体的微分方程。设 $t$ 是 Teichmüller 空间的复坐标（复结构参数），
$s_t\in H^0(\cM_J,L^{\otimes k}\otimes K^{1/2})$ 是一族截面。Hitchin 联络是
\begin{equation}
\nabla_{\partial_t}s=\partial_t s+\frac{1}{k+h^\vee}\left(\Theta_t\right)s, \tag{4.4}
\end{equation}
其中 $\Theta_t$ 是由 Beltrami 微分与模空间上的二阶基本形式构造的算符（具体形式依赖符号约定）。
这个联络的曲率是纯中心的：$[\nabla_{\partial_t},\nabla_{\partial_{\bar t}}]=\frac{\ii}{2}\omega_{\mathrm{WP}}$，
其中 $\omega_{\mathrm{WP}}$ 是 Weil--Petersson 形式。因此平行移动只差一个整体相位，
这正是“射影平坦”。

若把 $t$ 取成带标记点曲面的模参数，并把 $K^{1/2}$ 的贡献显式写出，(4.4) 就变成 KZ 方程 (6.5)。
这是第三、四讲的三维语言与第六讲的二维语言之间最直接的一座桥：Hitchin 联络的平坦性
（三维侧要求）等价于共形块的 Ward 恒等式（二维侧要求）。

\subsection{小结}

几何量子化的三步，在我们的情形里都有具体的答案：
\begin{enumerate}[leftmargin=*]
  \item 预量子化：行列式线丛，整性自动；
  \item 极化：Kähler 极化，量子态 $=H^0(\cM_J,L^{\otimes k})$；
  \item 半形式：把 $k$ 平移成 $k+h^\vee$。
\end{enumerate}
至此，三维理论的 Hilbert 空间有了一个完全具体的描述。下一讲把纽结放进来。

\subsection*{练习}
\begin{enumerate}[leftmargin=*]
  \item 验证 (4.1) 与 $k\in\ZZ$ 的等价性（提示：$\omega$ 的周期积分）。
  \item 对 $\RR^2$ 取 $P=\mathrm{span}\{\partial/\partial p\}$，写出满足 $\nabla_Xs=0$ 的截面，并说明它只依赖 $x$。
  \item 用 Riemann--Roch 计算 $g=1$、$SU(2)$、$L^{\otimes k}$ 的截面数，并说明半形式修正如何把它变成 $k+1$。
  \item 解释 (4.4) 中的分母为什么是 $k+h^\vee$ 而不是 $k$。
\end{enumerate}

\begin{quote}
\small
\textbf{本讲要点。} 量子态是 $\cM_J$ 上 $L^{\otimes k}$ 的全纯截面；$J$ 依赖只到射影平坦；
半形式修正把 $k$ 平移为 $k+h^\vee$。
\end{quote}

\section{纽结：Wilson 线与穿刺点}

\subsection{纽结从哪里进入}

第三、四讲处理的是一个没有纽结的三维流形：沿 $\Sigma$ 切开，得到一个干净的模空间。
现在要把一条 Wilson 线 $C$ 放进去。做法仍然是沿 $\Sigma$ 切开，但要选 $\Sigma$ 与 $C$ 横截相交；
于是 $\Sigma$ 上出现有限个交点 $P_1,\dots,P_n$，每个点带一个表示 $R_s$。

沿着 $C$ 的 Wilson 线 (2.5) 是“沿一条曲线的和乐”。在切开之后，这条曲线的两个端点
落在 $\Sigma$ 上，而中间部分在 $\Sigma$ 的一侧。如果我们只关心 $\Sigma$ 这一侧的量子态，
那么 Wilson 线的作用就是：在 $P_s$ 处插入一个“带表示 $R_s$ 的算符”。

\subsection{穿刺点上的高斯定律}

第三讲里高斯定律是 $F=0$。现在 $F$ 在穿刺点处不能为零：Wilson 线在那里是奇异的。
对 $A_0$ 变分，得到的约束变成
\begin{equation}
\frac{k}{8\pi}\,\varepsilon^{ij}F^a_{ij}(x)=\sum_s\delta^2(x-P_s)\,T^a_{(s)}, \tag{5.1}
\end{equation}
其中 $T^a_{(s)}$ 是表示 $R_s$ 的生成元。也就是说：曲率在 $P_s$ 处有一个“源”，
源的强度由表示决定。这正是电磁学里“点电荷给出电通量”的规范场版本。

\subsection{为什么不能经典地解}

(5.1) 看起来像一个方程，但它不是。原因有两条。

\textbf{第一，右端是算符。} $T^a_{(s)}$ 作用在表示空间 $R_s$ 上；不同的 $T^a$ 之间不对易。
因此 (5.1) 的右端不是一个数，而是一个算符值的分布。要“解”这个方程，意味着找到一个
数值联络 $A$ 使得两边相等，但右边的算符可以在 $R_s$ 上取不同的本征值，除非 $R_s$ 是一维的，
否则不存在一个同时满足所有分量的经典解。

\textbf{第二，$A$ 在 $P_s$ 处奇异。} 即使忽略算符问题，来源项也意味着 $A$ 在 $P_s$ 附近
不再是光滑的；平坦性只在 $\Sigma\setminus\{P_s\}$ 上成立。相空间因此是“带奇点的平坦联络模空间”，
而不是第三讲那个光滑的 $\cM$。

这两条合起来说明：穿刺点上的约束不是一个可以代入求解的方程，
而是一个必须在量子层面处理的“源条件”。

\subsection{把表示也量子化}

Witten 的处理办法是换一个视角：与其去解 (5.1)，不如把表示 $R$ 本身看成一个量子系统。
具体地，$R$ 是紧致群 $G$ 的不可约表示，而每一个不可约表示都可以实现为某个辛流形的几何量子化：
取余伴随轨道，或等价地取旗流形 $G/T$（$T$ 是极大环面）。$G/T$ 上有自然的 Kähler 结构，
对支配整权 $\lambda$，取辛形式 $\omega_\lambda$，其 de Rham 类就是 $\lambda$。
Borel--Weil--Bott 定理说：
\begin{equation}
H^0(G/T,\,\cL_\lambda)=R_\lambda, \tag{5.2}
\end{equation}
即 $G/T$ 上全纯线丛 $\cL_\lambda$ 的全纯截面空间，正是最高权为 $\lambda$ 的不可约表示。

注意这里 $\lambda$ 没有加 $\rho$：这是指标定理的直接结论。$\rho$ 的位移属于半形式修正
$K_{G/T}^{1/2}=\cL_{-\rho}$，不在 (5.2) 里。如果读者在自己的计算里看到 $\lambda+\rho$，
那多半是在用 Weyl 维数公式的写法，或者已经把半形式算进去了。

BWB 的证明思路值得知道，因为它解释了“为什么恰好是全纯截面”。$G/T$ 上的线丛 $\cL_\lambda$
的截面可以看成 $G$ 上的函数 $f$，满足 $f(gb)=\chi_\lambda(b)^{-1}f(g)$（$b\in B$，Borel 子群）。
按 Peter--Weyl 定理，这样的函数由最高权不超过 $\lambda$ 的表示张成；再要求全纯（即被
$\bar\partial$ 杀掉），就只剩下最高权恰为 $\lambda$ 的那个表示。于是 (5.2) 成立。
这条路线的关键是：不可约表示可以用“旗流形上的全纯截面”这个几何对象来实现，
而几何对象自然带辛结构，因而可以作为一个量子系统的相空间。

\subsection{抛物结构与扩大后的相空间}

把 (5.2) 代回，我们得到一个新的相空间：在 $\Sigma\setminus\{P_s\}$ 上取平坦 $G$-联络，
并要求在每个 $P_s$ 处，联络的结构群约化到 $T$（这样“源”就由 $\lambda_s$ 的整性给出）。
这个相空间记作 $\cM_{\{P_s,R_s\}}$。用丛的语言，这是在 $\Sigma$ 上取一个抛物 $G$-丛：
在 $P_s$ 处给定一个 Borel 子群（或一个滤过），平坦联络在这个滤过上取特定的形式。
抛物丛的模空间同样是 Kähler 的，只是多了边界（对应抛物权重的选择）。

在这个空间上，约束 (5.1) 变成了一个普通的经典约束（因为右端现在由 $G/T$ 的量子化提供，
不再是一个未定的算符），而辛形式也相应地扩大。把扩大后的相空间量子化，
得到的 Hilbert 空间记作 $\cH_{\Sigma;\{P_s,R_s\}}$。Witten 的论断是：
这个空间恰好是带有穿刺点 $P_s$、表示 $R_s$ 的共形块空间。
用第六讲的语言，它就是 WZW 模型的共形块。

\subsection{可积表示}

并非所有表示都给出非零的截面。只有“可积”的表示才有非零贡献；对 $SU(2)$，条件是最高的
自旋满足
\[
2j\le k.
\]
这个条件不是 Borel--Weil--Bott 的推论（BWB 对任意支配整权都成立），它来自仿射李代数：
level $k$ 的表示必须满足 $\lambda$ 在仿射 Weyl 群的基本区域里，而 $2j\le k$ 是它的
$SU(2)$ 特例。不可积的表示在量子化后给出零空间——用物理的语言说，这些表示“退耦”了。

\subsection{这是纲领，不是定理}

需要强调：上面这一段是 Witten 给出的\textbf{纲领}。原文用 "should" 描述它，
并明确说 "deferring a fuller treatment"——完整的构造留待后来。数学上把它严格化，
需要抛物丛、相交上同调、以及共形块理论；这些工作在 1990 年代之后才陆续完成。
本讲义把它作为理解的框架，而不是作为已证明的定理。

\subsection*{练习}
\begin{enumerate}[leftmargin=*]
  \item 说明 (5.1) 右端为什么是算符，并给出一个 $R$ 为一维表示时它可以经典求解的理由。
  \item 对 $G=SU(2)$、$R$ 为自旋 $j$ 的表示，写出 $\lambda$ 与 $j$ 的关系，并验证 BWB 给出的维数 $2j+1$。
  \item 为什么可积性条件是 $2j\le k$ 而不是 $2j\le k+2$？（提示：考虑半形式修正后哪个 $k$ 出现在条件里。）
\end{enumerate}

\begin{quote}
\small
\textbf{本讲要点。} Wilson 线在穿刺点处给出算符值的高斯定律 (5.1)；
Witten 用 $G/T$ 的几何量子化（BWB）把表示变成相空间的一部分，从而把约束经典化；
量子化后得到带穿刺点的共形块。可积性 $2j\le k$ 来自仿射代数。
\end{quote}

\section{二维一侧：WZW 模型与共形块}

\subsection{三维与二维的接口}

上一讲末尾，三维理论的 Hilbert 空间被写成了 $\cM_J$ 上全纯截面的空间。
这一讲要说明：同一个空间，也是二维共形场论的“共形块”空间。
两个理论在\textbf{态空间}层面是同一个东西；这不是说它们作为量子场论等价，
而是说同一个向量空间有两种描述。

二维一侧的好处是：那里的语言（表示、共形权、中心荷、融合规则）是纯代数的，
算起来比三维路径积分容易得多。所以这一讲的任务是把二维这套词汇搭起来，
让第七讲的每一次手术计算都有代数依据。读者如果只想记住一句话，那就是：
三维理论“知道”纽结，是因为它的态空间等于共形块空间，而共形块的融合规则就是纽结的拆接关系。

\subsection{WZW 作用量与 Wess--Zumino 项}

二维一侧的基本对象是 Wess--Zumino--Witten 模型：场 $g:\Sigma\to G$ 取值在紧致群中，
作用量由动能项与一个“拓扑”项相加：
\[
S_{\mathrm{WZW}}[g]
=\frac{k}{8\pi}\int_\Sigma \Tr\left(g^{-1}\partial g\;g^{-1}\bar\partial g\right)\dd^2z
+\frac{k}{12\pi}\int_B \Tr\left((g^{-1}\dd g)^3\right),
\]
其中 $\Sigma$ 是闭黎曼面，$B$ 是它的一个三维填充（$\partial B=\Sigma$）。
第一项是普通的 $\sigma$ 模型动能；第二项叫 Wess--Zumino 项。

它为什么必须写成三维积分？因为被积的 $\Tr(g^{-1}\dd g)^3$ 是一个\textbf{三形式}，
在二维流形 $\Sigma$ 上根本无从积分。用三维填充 $B$ 是它唯一的存在方式。
代价是 $B$ 不唯一：若 $B'$ 是另一个填充，$B\cup(-B')$ 是闭三维流形，
于是
\[
S_{\mathrm{WZW}}[g]\Big|_{B'}-S_{\mathrm{WZW}}[g]\Big|_{B}
=\frac{k}{12\pi}\int_{B\cup(-B')}\Tr\left((g^{-1}\dd g)^3\right)
=\frac{k}{12\pi}\cdot 24\pi^2\,n=2\pi k\,n,
\]
其中 $n$ 是映射 $g$ 的绕数。这一步用了 $\pi_3(G)\cong\ZZ$（紧致单连通单群）：
三形式的积分除以 $24\pi^2$ 给整数。因此作用量本身依赖填充，
但路径积分的权重 $\exp(\ii S)$ 良好定义，当且仅当
\[
k\in\ZZ .
\]
这正是第二讲里 $k$ 的整性来源，只是换了一件外衣：那边是大规范变换使
$\Delta L=2\pi k\,m$，这边是换填充使 $\Delta S=2\pi k\,n$；
两个整数 $m,n$ 都是同一个 $\pi_3(G)\cong\ZZ$ 的绕数。所以两处“$k$ 必须整”不是两个要求，
而是同一个要求。

\subsection{仿射李代数}

WZW 模型的手征流是
\[
J(z)=k\,\partial g\,g^{-1}\qquad(\text{反全纯流 } \bar J=k\,\bar\partial g\,g^{-1}),
\]
写成 Lie 代数分量 $J^a(z)$，它满足算子积展开
\begin{equation}
J^a(z)J^b(w)\sim\frac{k\,\delta^{ab}}{(z-w)^2}+\frac{\ii f^{abc}J^c(w)}{z-w}. \tag{6.1}
\end{equation}
把它按模式展开 $J^a(z)=\sum_{n\in\ZZ}J^a_n z^{-n-1}$，就得到无穷维代数
\[
[J^a_m,J^b_n]=\ii f^{abc}J^c_{m+n}+k\,m\,\delta^{ab}\,\delta_{m+n,0}.
\]
第一项是 Lie 代数的模式版本，第二项是\textbf{中心扩张}：
它是数（与 $J$ 无关），只在 $m+n=0$ 时出现，系数就是 level $k$。
这个代数记作 $\hat{\mathfrak g}_k$，叫仿射 Kac--Moody 代数。
没有第二项时，代数只是 $\mathfrak g$ 值函数的环路代数，没有最高权表示论可言；
中心扩张是整个二维一侧能“量子化”的关键。

\subsection{表示论、Weyl 群与 level 截断}

仿射代数的表示用最高权来分类。取一个有限维 $\mathfrak g$ 的不可约表示 $R_\lambda$
（最高权 $\lambda$），要求
\[
J^a_0|\lambda\rangle=t^a_\lambda|\lambda\rangle,
\qquad J^a_n|\lambda\rangle=0\quad(n>0).
\]
由 $J^a_{-n}$（$n>0$）作用生成的模叫 Verma 模。Verma 模一般可约：
某些组合会被零向量（null vector）打成零。把这些零向量商掉，得到不可约模
$\hat V_{k,\lambda}$。

何时 $\hat V_{k,\lambda}$ 是\textbf{酉}的（即内积正定）？答案是\textbf{可积}条件：
\[
(\lambda,\theta)\le k
\qquad\text{对所有正根 }\theta ,
\]
等价地 $(\lambda,\theta_{\max})\le k$，其中 $\theta_{\max}$ 是最大根。
满足条件的 $\lambda$ 叫 level $k$ 的可积最高权。几何上，这些 $\lambda$ 落在
仿射 Weyl 群
\[
W_a=W\ltimes kQ^\vee
\]
作用下的\textbf{基本 alcove} 内；Weyl 群并不直接作用在 $\lambda$ 上，
而是作用在平移后的 $\lambda+\rho$ 上（$\lambda\mapsto w(\lambda+\rho)-\rho$）。
\textbf{这正是 $\rho$ 与 $k+h^\vee$ 反复出现的原因}：所有公式都活在“平移过的”坐标里，
平移量是 $\rho$，而 $k$ 被 $h^\vee$ 补成 $k+h^\vee$ 才与 Weyl 群相容。

对 $SU(2)$，条件化简为 $2j\le k$，于是可积表示是
\[
j=0,\ \frac12,\ 1,\ \dots,\ \frac k2,
\]
一共 $k+1$ 个。对 $SU(N)$，可积表示的个数是
$\binom{k+N-1}{N-1}$（$SU(2)$ 时回到 $k+1$）。
这个“$k+1$”后面会反复出现：它是环面 Hilbert 空间的维数，是 $S$ 矩阵的阶数，
也是 Verlinde 公式里求和项的个数。

\subsection{Sugawara 构造与中心荷}

仿射代数上可以构造一个 Virasoro 代数，办法是把流“平方”：
\begin{equation}
T(z)=\frac{1}{2(k+h^\vee)}\sum_a :J^a(z)J^a(z):, \tag{6.2}
\end{equation}
其中 $:\ :$ 是正规序（把产生算子放左边），$h^\vee$ 是对偶 Coxeter 数。
这个构造叫 Sugawara 构造。它的合理性来自一个直接但繁琐的 Wick 收缩计算：
把 (6.1) 代进 $T(z)T(w)$，只保留四次极点，会发现
\begin{itemize}[leftmargin=*]
  \item 两次双收缩贡献的量级是 $k^2\dim G$；
  \item 一次双收缩加一次单收缩贡献的量级是 $k\,h^\vee\dim G$。
\end{itemize}
若只用双收缩，得到的“中心荷”是 $k^2\dim G/(2(k+h^\vee)^2)$，
它不满足 Virasoro 代数（$T$ 与自身的 OPE 不能闭合）。
补上单收缩项之后两项合并，恰好给出
\begin{equation}
c=\frac{k\dim G}{k+h^\vee}. \tag{6.3}
\end{equation}
分母里的 $k+h^\vee$ 不是凑出来的，它来自上述两类收缩的比值；
换句话说，$h^\vee$ 是流代数的“背景荷”，没有它 Virasoro 代数不闭合。

对 $SU(2)$，$c=3k/(k+2)$：$k=1$ 给 $c=1$，$k\to\infty$ 给 $c\to3$，
即三个自由玻色子——这与“大 $k$ 时理论趋于自由”的图像一致。
一般地 $k\to\infty$ 时 $c\to\dim G$。

表示 $R_\lambda$ 的共形权同样带 $k+h^\vee$：
\begin{equation}
h_\lambda=\frac{(\lambda,\lambda+2\rho)}{2(k+h^\vee)}. \tag{6.4}
\end{equation}
对 $SU(N)$ 的定义表示，$\lambda=\omega_1$，$(\omega_1,\omega_1)=(N-1)/N$、
$2(\omega_1,\rho)=N-1$，代入得
\[
h=\frac{N^2-1}{2N(N+k)} .
\]
这个公式在第七讲会以“Wilson 线的标架相位”再次出现，那里我们用它核对三点函数的指数。

\subsection{共形块与 Ward 恒等式}

在带穿刺点 $P_1,\dots,P_n$、每个点上带表示 $R_i$ 的黎曼面上，\textbf{共形块}是满足
Ward 恒等式的“手征关联函数”的线性空间。具体说，取场 $\phi_{\lambda_i}$ 的关联函数
\[
\Psi(z_1,\dots,z_n)=\left\langle\phi_{\lambda_1}(z_1)\cdots\phi_{\lambda_n}(z_n)\right\rangle ,
\]
要求它对手征坐标全纯，并且被能-动量张量的 Ward 恒等式约束：
\[
\left\langle T(z)\prod_i\phi_i(z_i)\right\rangle
=\sum_i\left[\frac{h_i}{(z-z_i)^2}+\frac{1}{z-z_i}\partial_{z_i}\right]
\left\langle\prod_i\phi_i\right\rangle .
\]
“解这个线性系统”就是“求共形块”。对有理共形场论，解空间是有限维的；
它的维数不依赖复结构（只依赖标记点的拓扑与表示），这本身就是一个非平凡的定理。

\subsection{KZ 方程从哪里来}

把 Sugawara 表达式 (6.2) 代进 (6.5) 的左边，用 (6.1) 收缩，就得到 $\Psi$ 的一阶方程。
这是本讲最值得亲手推一遍的计算，分四步。

\textbf{第一步：收缩。}$T(z)$ 与每个 $\phi_i$ 收缩一次，留下 $J^a(z)\phi_i(z_i)\sim
t^a_i\phi_i(z_i)/(z-z_i)$；对 $\sum_a:J^aJ^a:$ 来说，这就是单极点项的来源。

\textbf{第二步：分出双极点。}把 $T$ 写成 $z$ 的展开，$\frac{1}{(z-z_i)(z-z_j)}$ 型的项
在 $z\to z_i$ 时给出 $1/(z-z_i)^2$（取 $j=i$）与 $1/(z-z_i)$（取 $j\ne i$）两类。

\textbf{第三步：双极点与共形权比较。}双极点的系数应当是 $h_i$，比较后得到
$h_i=(\lambda_i,\lambda_i+2\rho)/2(k+h^\vee)$，即 (6.4)。这一步同时验证了 Sugawara 的归一化。

\textbf{第四步：单极点给出 KZ。}单极点的系数在两边的比较给出
\begin{equation}
\kappa\,\partial_{z_i}\Psi=\sum_{j\ne i}\frac{\Omega_{ij}}{z_i-z_j}\Psi,
\qquad \kappa=k+h^\vee,\quad \Omega_{ij}=\sum_a t^a_i t^a_j . \tag{6.5}
\end{equation}
这就是 Knizhnik--Zamolodchikov（KZ）方程。系数里的 $\kappa=k+h^\vee$ 与 (6.2) 的分母是同一个；
所以第四讲中“半形式把 $k$ 平移成 $k+h^\vee$”这件事，在二维侧就是 Sugawara 公式里的那个平移——
两边说的是同一件事。

\subsection{两点与四点：解长什么样}

两点情形最简单。KZ 方程 (6.5) 只剩一项，直接积分得
\[
\Psi(z_1,z_2)=(z_1-z_2)^{-2h_j},
\]
指数里的 $k$ 通过 $h_j$ 含在 $k+h^\vee$ 里。这正是二点函数
$\langle\phi_j(z_1)\phi_j(z_2)\rangle$ 的标准形式，说明 KZ 方程的解就是共形块。

三点情形：固定 $SL(2,\CC)$ 后三个点可以放到 $0,1,\infty$，KZ 方程退化为一组代数关系；
解空间至多一维，通常记为 $N_{ijk}\in\{0,1\}$（对 $SU(2)$）。
这正是下一小节要用的“三点共形块维数=融合系数”。

四点情形开始有意思。把 $SL(2,\CC)$ 固定到 $(0,z,1,\infty)$ 后，KZ 方程变成关于 $z$ 的
\textbf{二阶 Fuchs 型方程}，解空间是二维的（一般情形）。
两个独立解对应两种“中间表示”的交换：
$s$-道把 $1,2$ 先融合、$t$-道把 $2,3$ 先融合；两套解互为解析延拓，
其单值性（monodromy）就是辫群的表示。换言之，
\textbf{KZ 方程的解空间上带着一个平坦联络，它的单值性就是编织}。
三维一侧的“纽结”从这里进来：辫群的表示给出 Jones 多项式的 $R$ 矩阵。

有一个小例子可以手算：$SU(2)$、$k=1$、四个 $j=\frac12$。
因为 $2j\le k$ 只允许 $j=0,\frac12$，而 $\frac12\otimes\frac12=0$，
所以中间表示只能是 $0$，解空间一维，四点块被唯一确定（直到整体归一化）。
这解释了为什么 $k=1$ 的 Jones 多项式特别简单：只有一个通道。

\subsection{共形块的维数与融合系数}

对有理（rational）的共形场论，共形块空间有限维，维数由融合规则给出：
三点函数空间的维数是
\[
N_{ijk}=\dim\cH_{\{P_1,P_2,P_3;R_i,R_j,R_k\}},
\]
它满足融合代数
\[
R_i\otimes R_j=\bigoplus_k N_{ij}{}^kR_k ,
\qquad N_{ij}{}^k=N_{ijk}\ \text{（用对偶度规升降指标）}.
\]
一般地，$n$ 点共形块的维数是把 $n$ 个表示按融合规则配成不变量的方案数：
\[
\dim\cH_{S^2;R_1,\dots,R_n}
=\sum_{i_1,\dots,i_{n-3}}
N_{i_1 i_2}{}^{j_1}N_{j_1 i_3}{}^{j_2}\cdots .
\]
这个“数方案”的图像在第七讲会变成“手术时插入完备态然后求和”，
两者是同一件事的两种说法。

\subsection{特征标与模不变性}

共形块在环面上求迹，得到特征标
\begin{equation}
\chi_\lambda(\tau)=\Tr_{\hat V_{k,\lambda}}\left(q^{L_0-c/24}\right),
\qquad q=e^{2\pi\ii\tau}. \tag{6.6}
\end{equation}
$q^{L_0-c/24}$ 里的 $c/24$ 是柱面上的真空能：$L_0$ 在真空上取值 $-c/24$ 而不是零，
这一点在第七讲计算标架反常时会再遇到。

模群 $\mathrm{SL}(2,\ZZ)$ 由 $\tau\to\tau+1$ 与 $\tau\to-1/\tau$ 生成，它把特征标混合起来：
\[
\chi_\lambda(\tau+1)=e^{2\pi\ii(h_\lambda-c/24)}\chi_\lambda(\tau),
\qquad
\chi_\lambda(-1/\tau)=\sum_\mu S_{\lambda\mu}\,\chi_\mu(\tau).
\]
$T$ 变换给对角相位，$S$ 变换给混合矩阵。对 $SU(2)$、level $k$，$S$ 有闭式
\begin{equation}
S_{jj'}=\sqrt{\frac{2}{k+2}}\,\sin\frac{\pi(2j+1)(2j'+1)}{k+2},
\qquad j,j'=0,\frac12,\dots,\frac k2 . \tag{6.7}
\end{equation}
模不变性（配分函数在模变换下不变）给出一个强约束：$S$ 必须是对称酉矩阵，
并且满足模群关系。用 Witten 的 $T=\mathrm{diag}\big(e^{2\pi\ii(h_j-c/24)}\big)$ 时，
\begin{equation}
S^2=(ST)^3=1 ; \tag{6.8}
\end{equation}
这一点可以手算验证。取 $k=2$，由 (6.7) 得
\[
S=\begin{pmatrix}
\frac12 & \frac{\sqrt2}{2} & \frac12\\[3pt]
\frac{\sqrt2}{2} & 0 & -\frac{\sqrt2}{2}\\[3pt]
\frac12 & -\frac{\sqrt2}{2} & \frac12
\end{pmatrix}.
\]
逐项相乘：$S^2$ 的 $(0,0)$ 元是 $\frac14+\frac12+\frac14=1$；
$(0,1)$ 元是 $\frac{\sqrt2}{4}+0-\frac{\sqrt2}{4}=0$；
$(1,1)$ 元是 $\frac12+0+\frac12=1$；
$(1,2)$ 元是 $\frac{\sqrt2}{4}+0-\frac{\sqrt2}{4}=0$。
其余同理，故 $S^2=1$。这个验算值得每个人自己做一遍：
它同时检查了 (6.7) 的归一化因子 $\sqrt{2/(k+2)}$ 与“可积表示个数 $=k+1$”两件事。

有了 $S$，融合规则就被它对角化（Verlinde）：
\begin{equation}
N_{ijk}=\sum_{l}\frac{S_{il}\,S_{jl}\,S_{kl}}{S_{0l}} . \tag{6.9}
\end{equation}
这一步的推导是第七讲的主题；在二维一侧，它的意思是“$S$ 矩阵的列是融合代数的特征向量”。
读者可以把 (6.9) 当作一句承诺：只要知道 $S$，就知道所有融合规则，
于是所有共形块维数都可以算。

\subsection{三个最小例子的数据表}

把 $k=1,2,3$ 的数据列成表，第七讲计算 $S^3$ 配分函数时会直接用到。

\begin{center}
\small
\begin{tabular}{@{}clcccc@{}}
\toprule
$k$ & 可积 $j$ & 表示个数 & $c=3k/(k+2)$ & $h_j=\dfrac{j(j+1)}{k+2}$ & $S$ 矩阵 \\
\midrule
1 & $0,\frac12$ & 2 & $1$ & $0,\ \frac14$ & $\frac{1}{\sqrt2}\begin{pmatrix}1&1\\1&-1\end{pmatrix}$ \\
2 & $0,\frac12,1$ & 3 & $\frac32$ & $0,\ \frac{3}{16},\ \frac12$ & $\frac{\sqrt2}{2}\begin{pmatrix}\frac{1}{\sqrt2}&1&\frac{1}{\sqrt2}\\1&0&-1\\ \frac{1}{\sqrt2}&-1&\frac{1}{\sqrt2}\end{pmatrix}$ \\
3 & $0,\frac12,1,\frac32$ & 4 & $\frac95$ & $0,\ \frac{3}{20},\ \frac25,\ \frac34$ & 由 (6.7) 给出 \\
\bottomrule
\end{tabular}
\end{center}

顺带给出三个检验，读者可以自己核：
\begin{enumerate}[leftmargin=*]
  \item $\sum_jS_{0j}^{\,2}=1$ 对每个 $k$ 都成立（$S$ 酉且实对称，第一行是单位向量）。
  \item $h_j$ 的公式在 $k\to\infty$ 时趋于 $j(j+1)/k$，趋于经典值；
        $c$ 在 $k\to\infty$ 时趋于 $3=\dim SU(2)$。
  \item 表中 $h_j$ 的分母是 $k+2$，不是 $k$——凡是写成 $j(j+1)/k$ 的公式都把裸 level
        与有效 level $k+h^\vee$ 混了。
\end{enumerate}

关于这些理论的“身份”，两点值得注意。
$k=1$ 时 $c=1$，$SU(2)_1$ 等价于半径 $R=\sqrt2$ 的紧致自由玻色子（自对偶点）。
$k=2$ 时 $c=3/2$，与三份 Ising 模型有相同的中心荷，但两者初级场谱不同
（$SU(2)_2$ 的共形权是 $0,\frac{3}{16},\frac12$，Ising 的是 $0,\frac{1}{16},\frac12$），
所以不是直和；把 $SU(2)_2$ 的 $U(1)$ 子代数商掉，剩下的正是 Ising 模型
（$Z_2$ 仲费米子）。$k=3$ 时 $c=9/5>1$，它不属于最小模型（最小模型的 $c<1$），
这一点常被误记。

\subsection{同一个空间}

现在回到第四讲：$\cH_\Sigma=H^0(\cM_J,L^{\otimes k})$。
Witten 的论断是：这个空间与 WZW 模型的共形块空间\textbf{是同一个向量空间}。
两者都是“带标记点曲面的共形块”，都带同一个射影平坦丛，
而且在标记点处的插入由同一个表示给出。

这个等同可以这样理解：$\cM_J$ 上 Hitchin 联络的平坦性，正好等价于共形块满足 Ward 恒等式；
射影平坦性（共形 anomaly）对应中心荷 (6.3)；而穿刺点上的抛物结构对应 WZW 的插入。
换句话说，第四讲里“$\cH_\Sigma$ 射影平坦”这句话的二维翻译，就是“存在共形块”。

\subsection{这个等同是什么级别}

需要精确地说：这是一个\textbf{对象层面}的等同——两个向量空间同构，
并且同构与射影平坦丛相容。它\textbf{不是}两个理论作为量子场论的等价：
三维理论没有边界上的动力学，二维理论有；三维理论没有引力子，
二维理论的共形块也不描述动力学。

更强的“函子性”（把曲面的缝合对应到空间的张量积）在原文里只被提到一句
“It is obvious at this point”，严格化是后来的工作（Segal 的模函子公理、
Moore--Seiberg 的缝合证明）。所以我们用“同一个向量空间”而不用“同一个理论”，
是有意的措辞。

\subsection{从 WZW 到极小模型：陪集构造}

WZW 模型是“母模型”，很多有理共形场论可以从它商出来。设 $H\subset G$ 是一个子群，
把 $G$-WZW 模型按 $H$ 的规范对称性商掉，得到陪集（coset）模型 $G/H$。
它的中心荷是
\[
c_{G/H}=c_G-c_H=\frac{k_G\dim G}{k_G+h^\vee_G}-\frac{k_H\dim H}{k_H+h^\vee_H},
\]
其中 $k_H$ 由嵌入的“水平”决定（对嵌入 $H\subset G$，水平按 Dynkin 指数的比例缩放）。

最重要的例子是 $SU(2)_k/U(1)$。它的中心荷是 $c=1-\frac{6}{(k+2)(k+3)}$，
这正是极小模型 $M(k+1,k+2)$ 的中心荷。也就是说，极小模型（包括 Ising 的 $c=\frac12$、
三态 Potts 的 $c=\frac45$）都可以写成 $SU(2)$ WZW 模型的陪集。

这个构造对本文的意义是：它说明 WZW 模型不是“一个模型”，而是一个“模型工厂”。
本文用 WZW 来解释 Chern--Simons，但 WZW 本身又生成了一大批二维模型；
这条链条（三维 CS → 二维 WZW → 二维极小模型）是 1980 年代后期共形场论的核心图景。

\subsection{与自由场的关系：玻色化}

WZW 模型不是凭空出现的，它有自由场实现。最简单的例子是 $SU(2)_1$：
它的中心荷 $c=1$，与一个自由玻色子相同。事实上，$SU(2)_1$ 的 WZW 模型等价于
一个半径为 $R=\sqrt2$ 的紧致玻色子；两个可积表示 $j=0,\frac12$ 对应玻色子的两种
“缠绕/动量”扇区，共形权分别是 $0$ 与 $\frac14$。

这个等价关系叫玻色化（bosonization）。它有两个用处。第一，它给了 WZW 模型一个
可以显式计算的模型：自由场的关联函数是已知的，所以 $SU(2)_1$ 的共形块可以手算。
第二，它解释了为什么 $SU(2)_1$ 的共形块维数与 theta 函数有关：自由玻色子的态空间
就是 theta 函数空间。

对一般的 $G$ 与 $k$，情况更复杂，但有一个一般的构造：WZW 模型可以写成
“$G$-值场 + 一个 $G$ 的 Wess--Zumino 项”，在 $k=1$ 时可以用自由场表示。
更高 level 需要引入“轨道流形”或“陪集构造”，那是另一套技术。

对本文的读者来说，记住两点就够了：第一，WZW 不是抽象的公理系统，它有具体的自由场实现；
第二，$k=1$ 与自由场的关系，是理解 $c=1$ 与 $c=k\dim G/(k+h^\vee)$ 的钥匙。

\subsection{小结}

二维一侧给出的语言是：仿射代数与 level、可积表示与截断、Sugawara 与中心荷、
共形块与 KZ 方程、融合规则与 $S$ 矩阵。它们的逻辑关系是：
\begin{itemize}[leftmargin=*]
  \item level $k$ 决定代数 $\hat{\mathfrak g}_k$，也决定可积表示只有有限个；
  \item Sugawara 构造把流代数变成 Virasoro，中心荷 $c=k\dim G/(k+h^\vee)$；
  \item 共形块是 Ward 恒等式的解，等价于 KZ 方程 (6.5) 的全纯解；
  \item 共形块维数就是融合系数 $N_{ijk}$，而 $N_{ijk}$ 由 $S$ 矩阵经 (6.9) 给出。
\end{itemize}
第七讲会把这套代数语言与三维手术对接：每一次手术对应一次融合，
每一个链环分量对应一个表示，最后得到的数就是 Jones 多项式在 $q=e^{2\pi\ii/(k+N)}$ 处的值。

\subsection*{练习}
\begin{enumerate}[leftmargin=*]
  \item 由 (6.2) 与 (6.1) 验证 $c=k\dim G/(k+h^\vee)$ 对 $SU(2)$ 给出 $3k/(k+2)$。
  \item 计算 $SU(2)$、$k=1$ 时的 $h_0,h_{1/2}$ 与 $c$，并验证 $\sum_j S_{0j}^2=1$。
  \item 证明两点 KZ 方程的解是 $(z_1-z_2)^{-2h_j}$。
  \item 验证 $SU(2)$、$k=2$ 的 $S$ 矩阵满足 $S^2=1$，并用 (6.9) 算出 $N_{1,1}{}^{0}$ 与 $N_{1,1}{}^{1}$。
  \item 说明为什么 Wess--Zumino 项不能写成 $\Sigma$ 上的二维积分；
        再说明为什么换填充改变 $S$ 恰好是 $2\pi k$ 的整数倍。
  \item 用 (6.7) 算出 $k=3$ 的 $4\times4$ $S$ 矩阵的第一行，并验证其平方和为 1。
  \item 判断：$SU(2)_2$ 与三份 Ising 模型是同一个共形场论吗？给出理由。
\end{enumerate}

\begin{quote}
\small
\textbf{本讲要点。} WZW 的 level $k$ 给出仿射代数 $\hat{\mathfrak g}_k$ 与有限多个可积表示；
Sugawara 构造给中心荷 $c=k\dim G/(k+h^\vee)$ 与共形权 (6.4)；
共形块满足 KZ 方程 (6.5)，系数 $\kappa=k+h^\vee$；
共形块维数等于融合系数，融合系数由 $S$ 矩阵经 (6.9) 给出；
三维 Hilbert 空间与共形块空间在对象层面等同。
\end{quote}

\section{算出来：手术、$S$ 矩阵与 Verlinde 公式}

\subsection{要算什么}

到现在为止，我们有了一个理论：三维 Chern--Simons 理论，量子态是全纯截面。
但“理论”必须能算出数。这一讲算两类东西：
\begin{itemize}[leftmargin=*]
  \item 闭三维流形的配分函数 $Z(M)$；
  \item 带穿刺点曲面的 Hilbert 空间维数，即融合系数 $N_{ij}{}^k$。
\end{itemize}
算完之后，把纽结放进去，就得到 Jones 多项式。

\subsection{黎曼球上的三点函数}

先看最简单的曲面：$S^2$ 上三个标记点，表示分别为 $R_i,R_j,R_k$。对应的 Hilbert 空间维数记为
\begin{equation}
N_{ijk}=\dim\cH_{S^2;\{i,j,k\}}. \tag{7.1}
\end{equation}
这个数就是融合系数。$N_{ijk}\ne0$ 当且仅当 $R_k$ 出现在 $R_i\otimes R_j$ 的分解里。
在最简单的情形（大 $k$ 极限），$N_{ijk}$ 就是经典张量积的重数；有限 $k$ 时，它是它的一个子集
（只有可积表示才计入）。

\subsection{环面、$S$ 与 $T$}

亏格 $1$ 的情形给出另一个基本数据。环面 $\Sigma=T^2$ 上的 Hilbert 空间维数是 $k+1$（$SU(2)$ 时），
有一组自然基 $|j\rangle$，$j=0,\frac12,\dots,\frac k2$。模群 $SL(2,\ZZ)$ 作用在这个空间上，
由两个生成元给出：
\begin{equation}
T=\mathrm{diag}\left(e^{2\pi\ii(h_j-c/24)}\right),\qquad
S_{jj'}=\sqrt{\frac{2}{k+2}}\,\sin\frac{\pi(2j+1)(2j'+1)}{k+2}. \tag{7.2}
\end{equation}
$S$ 与 $T$ 满足 $S^2=(ST)^3=1$（在 Witten 的 $T$ 约定下；若把 $T$ 换成拓扑自旋
$\mathrm{diag}(e^{2\pi\ii h_j})$，则 $(ST)^3=e^{2\pi\ii c/8}$，差的正是标架反常）。

$S$ 的物理意义是“把环面沿一个圈切开再重新粘合”的变换。它也可以由共形块直接算出：
$S_{jj'}$ 是环面上的两个基态之间的内积，在模变换下交换了两个圈。

\subsection{手术}

任意一个闭三维流形都可以由 $S^3$ 沿一个带标架的链环做手术得到：挖掉每个链环分量的管状邻域
（实心环），再按标架换一种方式粘回去。每一步手术在量子理论里对应一个操作：
在切开处插入一个完备态，粘合时做求和。

具体地，沿一个纽结做手术时，插入的态是 $\sum_\lambda S_{0\lambda}\,|\lambda\rangle$ 型的组合，
其中 $S_{0\lambda}$ 是真空与 $\lambda$ 之间的 $S$ 矩阵元。于是
\begin{equation}
Z(M)=\sum_{\lambda_1,\dots,\lambda_r}
\left(\prod_{i}S_{0\lambda_i}\right)\;Z\!\left(S^3;\,\text{带标架的链环标号 }\lambda_i\right). \tag{7.3}
\end{equation}
这个公式把任意流形的配分函数化归为 $S^3$ 上带标号链环的配分函数，而后者可以用辫群与
共形块算出来。

\subsection{Verlinde 公式的推导}

现在考虑一个具体的比较：把 $S^2\times S^1$ 沿一个圈做手术，可以得到 $S^3$。
一方面，$S^2\times S^1$ 的 Hilbert 空间已知（就是 $S^2$ 的态空间）；另一方面，
手术之后应该得到 $S^3$ 的配分函数。两边相等，给出一个关于 $S$ 与 $N$ 的恒等式。

Witten 的推导（原文 (4.51)）给出
\begin{equation}
\sum_k S_i{}^k\,N_{kl}{}^m=\frac{S_{il}\,S_i{}^m}{S_{0i}}. \tag{7.4}
\end{equation}
这正是“$S$ 对角化融合规则”的数学表述。令 $v_i$ 为融合代数的基，$v_iv_j=\sum_k N_{ij}{}^kv_k$；
则在新基
\[
w_i=S_{0i}\sum_m S_i{}^m\,v_m
\]
下，乘法变成
\begin{equation}
w_iw_j=\delta_{ij}\,w_j, \tag{7.5}
\end{equation}
即 $w_i$ 是幂等元。于是所有融合系数都可以由 $S$ 算出：
\[
N_{ij}{}^k=\sum_l\frac{S_{il}S_{jl}S_{kl}}{S_{0l}}.
\]
这就是 Verlinde 公式。

\subsection{这个推导用了什么}

值得注意的是，上面这条推导链条只用到了三样东西：
\begin{enumerate}[leftmargin=*]
  \item 手术公式 (7.3)；
  \item $N_{ijk}$ 的定义 (7.1)；
  \item 态空间基的正交归一性。
\end{enumerate}
它\textbf{没有用到共形块的动力学}（KZ 方程、Ward 恒等式）。换句话说，Verlinde 公式的
“物理推导”实际上是拓扑场论自洽性的推论；把它变成纯数学定理，需要证明共形块空间确实满足
这些自洽性（TUY 1989、Faltings 1994 做的是这件事）。

\subsection{标架反常}

配分函数还依赖标架。改变标架时，
\begin{equation}
Z\longmapsto Z\cdot\exp\left(2\pi\ii\,\frac{n-m}{24}\,c\right), \tag{7.6}
\end{equation}
其中 $n-m$ 是标架改变对应的整数，$c$ 是 (6.3) 的中心荷。这个相位不影响 $|Z|$，
但影响带相位的量（例如 Jones 多项式在 $q$ 的单位根处的取值）。第八讲会再讨论它。

\subsection{环面配分函数与模不变性}

配分函数不只在闭三维流形上有定义；把时空取成 $\Sigma\times S^1$，得到的是一个二维配分函数
（对 $\Sigma$ 上的场求和），它作为 $S^1$ 的周长（模参数 $\tau$）的函数，满足模不变性。
对 $\Sigma=T^2$ 且无 Wilson 线的情形，
\[
Z_{T^2\times S^1}(\tau)=\sum_{j} |\chi_j(\tau)|^2 ,
\]
其中 $\chi_j$ 是可积表示的特征标，求和遍历可积表示。这个量在 $\tau\to-1/\tau$ 下不变，
因为 $S$ 矩阵是酉的、且特征标按 $S$ 变换。

这个公式是“模不变性”的最简形式，也是 Verlinde 公式的来源之一：
如果在环面上插入 Wilson 线（即标记点），配分函数变成 $\sum_j \chi_j\bar\chi_{j'}$ 的组合，
而融合系数正是这些组合的系数。换句话说，Verlinde 公式可以从“环面配分函数的模不变性”
推出来，这正是 Verlinde 原始论证的思路。

本文的推导走的是另一条路（手术 + 三维态空间），两条路给出同一个公式。
这种“多路验证”是有理共形场论的典型特征：一个公式常常有物理、代数、几何三种证明。

\subsection{一个可核对的数}

取 $G=SU(2)$、$k=2$。此时 $S_{00}=\sqrt{2/4}\sin(\pi/4)=1/2$，而 $S^3$ 的配分函数正是 $S_{00}$，
所以
\[
Z(S^3)=\tfrac12\qquad(k=2).
\]
一般地 $Z(S^3)=\sqrt{2/(k+2)}\sin(\pi/(k+2))$；$k=1$ 时它是 $1/\sqrt2$。
这个数可以用 (7.4) 与 $N$ 的定义独立核对，读者不妨自己算一遍。它是检验整套机器是否自洽的
一个最小例子。

\subsection{$S$ 矩阵从哪来}

(7.2) 的公式看起来像从天而降。其实它有一个干净的来源：环面态空间上的模变换。

取环面 $\Sigma=T^2$，记它的两个圈为 $a,b$。环面上的共形块（也就是三维理论在该曲面上的态）
由不可积表示 $|j\rangle$ 标号。模群 $SL(2,\ZZ)$ 由
\[
S=\begin{pmatrix}0&-1\\1&0\end{pmatrix},\qquad
T=\begin{pmatrix}1&1\\0&1\end{pmatrix}
\]
生成，它作用在态空间上。物理上，$T$ 是“沿 $a$ 圈做 Dehn 扭转”，
而 $S$ 是“交换 $a$ 与 $b$”。用共形场论的语言，$T$ 作用是对角矩阵
$\mathrm{diag}(e^{2\pi\ii(h_j-c/24)})$，因为扭转一圈给出共形权与中心荷的相位；
$S$ 作用则给出 (7.2) 的矩阵元。

(7.2) 的显式形式可以用 Weyl 特征标公式推出：$S_{jj'}$ 是特征标在模变换下的核，
对 $SU(2)$ 它就是 $\sqrt{2/(k+2)}\sin(\pi(2j+1)(2j'+1)/(k+2))$。
读者可以验证 $S^2=1$ 与 $(ST)^3=1$（在 Witten 的 $T$ 约定下）；这两个等式不是巧合，
它们就是模群关系的量子化版本。若把 $T$ 换成拓扑自旋，则 $(ST)^3$ 多出 $e^{2\pi\ii c/8}$，
这正是 (7.6) 的标架反常。

\subsection{辫群与 Jones 多项式}

最后一步是把链环写成辫子的闭合。任何定向链环都可以表示为一个辫子 $\beta\in B_n$ 的闭合
$\hat\beta$。在共形场论里，辫群作用在 $n$ 个标记点的共形块空间上：$B_n$ 的生成元
$\sigma_i$ 对应“把第 $i$ 与第 $i+1$ 个标记点交换”，而 KZ 方程给出的单值性正是这个作用。
闭合辫子对应于取迹：
\[
Z(S^3;\hat\beta)=\Tr_{\cH}\left(\rho(\beta)\right).
\]
对 $SU(2)$、定义表示，这个表示可以约化到 Temperley--Lieb 代数，其迹给出 Jones 多项式。
这就是 Jones 多项式“来自三维”的确切含义：它是共形块空间上辫群表示的一个迹。

\subsection{手术公式的推导}

把 (7.3) 写得更具体一些。设 $M$ 由 $S^3$ 沿链环 $L=L_1\cup\cdots\cup L_r$ 做手术得到，
标架为 $f_i$。手术的每一步可以分解为：
\begin{enumerate}[leftmargin=*]
  \item 挖去 $L_i$ 的管状邻域 $N_i\cong D^2\times S^1$，剩下一个边界为 $\coprod_i T^2_i$ 的流形 $M_0$；
  \item 在 $M_0$ 的每个边界环面上，态空间有一组基 $|\lambda_i\rangle$（由可积表示标号）；
  \item 实心环 $D^2\times S^1$ 的态是 $\sum_\lambda S_{0\lambda}|\lambda\rangle$
        （这是“真空在手术下的展开”，$S_{0\lambda}$ 是真空与 $\lambda$ 的 $S$ 矩阵元）；
  \item 粘合时做内积，得到求和。
\end{enumerate}
于是
\[
Z(M)=\sum_{\lambda_1,\dots,\lambda_r}
\left(\prod_i S_{0\lambda_i}\right)\,
Z(S^3;L\text{ 标号 }\lambda_1,\dots,\lambda_r).
\]
这个公式的严格版本需要证明：手术不依赖具体的链环表示（Kirby  calculus）。
Kirby 定理给出链环图之间的三种基本移动，而 $S$ 矩阵与融合系数满足的恒等式
（$S^2=(ST)^3=1$、Verlinde 公式）恰好保证 $Z(M)$ 在这些移动下不变。
这就是“为什么手术公式是自洽的”的答案：自洽性等价于模群关系。

\subsection{Heegaard 分裂与 WRT 不变量}

手术公式有一个等价的表述，叫 Heegaard 分裂。任何闭三维流形都可以写成两个柄体
沿一个曲面 $\Sigma$ 粘合：$M=M_1\cup_\Sigma M_2$，其中 $M_1,M_2$ 各是一个柄体。
在 Chern--Simons 理论里，每个柄体给 $\Sigma$ 上一个态（向量），粘合就是内积：
\[
Z(M)=\langle \psi_1,\psi_2\rangle .
\]
这个表述比手术公式更几何：它把三维流形的不变量化归为曲面上的态空间与两个向量。
而曲面上的态空间，正是第六讲的共形块空间。

Reshetikhin--Turaev 的构造可以看作这个表述的组合化：用一个模张量范畴（或量子群在单位根处的
表示范畴）替换“共形块”，用范畴的缝合运算替换“内积”，得到对链环与三维流形的组合不变量。
这个构造的优点是完全组合、可计算；缺点是它把“为什么”丢掉了——范畴从哪里来、
为什么与 Chern--Simons 一致，需要 Witten 的路径积分来解释。

对读者来说，两条路都值得知道：手术公式告诉你“怎么算”，Heegaard 分裂告诉你“为什么这样算”。
本文的贡献是用后者解释前者：三维理论在曲面上的态空间就是共形块，所以手术公式不是技巧，
而是理论的结构。

\subsection{一般群与 $S$ 的公式}

对一般的紧致单群 $G$ 与 level $k$，$S$ 矩阵由 Weyl 特征标公式给出：
\begin{equation}
S_{\lambda\mu}=\frac{1}{|W|}\sum_{w\in W}\varepsilon(w)\,
\exp\left(-\frac{2\pi\ii}{k+h^\vee}\langle w(\lambda+\rho),\mu+\rho\rangle\right), \tag{7.7}
\end{equation}
其中 $W$ 是 Weyl 群，$\varepsilon(w)$ 是符号，求和遍历 Weyl 群。
对 $SU(2)$，$W=\{1,-1\}$，把 $\lambda=2j$ 代入就回到 (7.2) 的正弦形式。
对 $SU(N)$，$S$ 矩阵可以用 Schur 多项式与单位根处的取值写出；
这些公式在 Verlinde 的原始论文与后来的综述里都有。

(7.7) 的一个重要推论是：$S$ 矩阵只依赖 level $k$ 与表示，不依赖曲面的亏格。
亏格只影响态空间的维数（由 Verlinde 公式给出），不影响融合规则。
这就是“有理共形场论”的结构：数据是有限维的、亏格无关的。

\subsection{三叶结：把数算出来}

最后用一个具体例子把整套机器串起来。三叶结是最简单的非平凡纽结，
它的 Jones 多项式是
\[
J(\text{三叶结})=-t^{-4}+t^{-3}+t^{-1}.
\]
在 Chern--Simons 的框架里，三叶结是 $2$ 股辫子 $\sigma_1^3$ 的闭合。
$SU(2)$、定义表示的辫群表示可以约化到 Temperley--Lieb 代数 $TL_n$，
其生成元 $e_i$ 满足
\[
e_i^2=\delta\,e_i,\qquad e_ie_{i\pm1}e_i=e_i,\qquad e_ie_j=e_je_i\ (|i-j|\ge2),
\]
其中 $\delta=-q-q^{-1}$，$q=e^{2\pi\ii/(k+2)}$。
辫子生成元 $\sigma_i$ 可以写成 $q^{1/2}(1-e_i)$ 的倍数；闭合取迹（Markov 迹）后，
得到的量就是 Jones 多项式在 $q$ 处的值。

取 $k=1$（此时 $q=e^{2\pi\ii/3}$），$TL_2$ 只有两个基，$e_1$ 是一个 $2\times2$ 矩阵；
$\sigma_1^3$ 的迹可以直接算出，得到 $J$ 在 $q$ 处的值。
把 $q$ 看作变量并解析延拓，就得到上面的多项式。

这个计算的意义在于：它不是“验证”，而是“定义”。在 Chern--Simons 框架里，
Jones 多项式的变量 $t$ 就是 $q^2$，它的单位根取值对应有限的 $k$；
多项式本身是这些取值的插值。这就是第一讲所说的“三维定义”的具体含义。

\subsection{手算 $k=2$ 的 $S$ 矩阵}

$k=2$ 时 $SU(2)$ 的可积表示是 $j=0,\frac12,1$，共三个。用 (7.2)：
\[
S_{jj'}=\sqrt{\frac{2}{4}}\,\sin\frac{\pi(2j+1)(2j'+1)}{4}
=\frac{1}{\sqrt2}\sin\frac{\pi(2j+1)(2j'+1)}{4}.
\]
逐项代入：
\[
S=\frac{1}{\sqrt2}
\begin{pmatrix}
\sin\frac{\pi}{4} & \sin\frac{2\pi}{4} & \sin\frac{3\pi}{4}\\[2pt]
\sin\frac{2\pi}{4} & \sin\frac{4\pi}{4} & \sin\frac{6\pi}{4}\\[2pt]
\sin\frac{3\pi}{4} & \sin\frac{6\pi}{4} & \sin\frac{9\pi}{4}
\end{pmatrix}
=
\frac12
\begin{pmatrix}
1 & \sqrt2 & 1\\
\sqrt2 & 0 & -\sqrt2\\
1 & -\sqrt2 & 1
\end{pmatrix}.
\]
于是 $S_{00}=\frac12$，与前面 $Z(S^3)=S_{00}=\frac12$ 一致；$S^2=1$ 可以逐项验证
（注意 $\sqrt2\cdot\sqrt2=2$），$(ST)^3=1$ 需要先写出 $T$，留给读者。

这个矩阵还有一个用途：用它算融合系数。例如
\[
N_{1/2,1/2,1}=\sum_{l=0,\frac12,1}\frac{S_{1/2,l}^2\,S_{1,l}}{S_{0,l}}
=\frac{\tfrac12\cdot\tfrac12}{\tfrac12}+\frac{\tfrac12\cdot\tfrac12}{\tfrac12}=1,
\]
第一项来自 $l=0$（$S_{1/2,0}^2=\tfrac12$，$S_{1,0}=\tfrac12$，$S_{0,0}=\tfrac12$），
第二项来自 $l=1$（$S_{1/2,1}^2=\tfrac12$，$S_{1,1}=\tfrac12$），$l=\frac12$ 项为零。
同理 $N_{1/2,1/2,0}=1$。于是 $1/2\otimes1/2=0\oplus1$，
与第六讲的融合规则一致。

\subsection*{练习}
\begin{enumerate}[leftmargin=*]
  \item 用 (7.2) 验证 $S^2=1$（$SU(2)$、$k=1,2$）。
  \item 计算 $Z(S^3)$ 在 $k=1,2,3$ 时的值，并验证它们等于 $S_{00}$。
  \item 用 (7.4) 与 $k=1$ 的 $S$ 矩阵验证 $N_{1/2,1/2,0}=1$、$N_{1/2,1/2,1/2}=0$。
\end{enumerate}

\begin{quote}
\small
\textbf{本讲要点。} 手术公式把 $Z(M)$ 化归为 $S^3$ 上的链环；$S$ 与 $T$ 由环面态空间给出 (7.2)；
(7.4) 是 Verlinde 公式的等价形式；标架反常 (7.6) 带 $c/24$；Jones 多项式是辫群表示在共形块上的迹。
\end{quote}

\section{边界与延伸：哪些是原文，哪些是后话}

\subsection{原文留下的空白}

读完前七讲，读者可能会觉得一切都很顺：作用量、相空间、量子化、共形块、计算。
但原文自己也清楚，有几处它只给了纲领而没有证明。把它们列出来，比记住结论更重要。

\textbf{第一，函子性。} 原文在讨论缝合时只写了一句 "It is obvious at this point"，
意思是：把曲面沿边界粘起来，对应的 Hilbert 空间应当做张量积。这个陈述把理论提升为
一个“模函子”（modular functor），但原文没有给出证明。严格化是后来的事（Segal 的框架、
以及大量关于共形块的手术公式的工作）。

\textbf{第二，标架反常。} (7.6) 的相位在一般手术下如何追踪，原文只处理了简单情形，
把一般情形留给了后续。这一项后来成为 Reshetikhin--Turaev 不变量里“缎带结构”的来源。

\textbf{第三，半形式修正。} 第四讲说过，$k\to k+h^\vee$ 的平移原文没有在正文里做；
它是后世（以及物理上的标架反常）补上的。

\textbf{第四，量子群。} 原文在讨论穿刺点约束时提到“量子群”，但没有构造。
后来的发展表明，可积表示范畴与量子群 $U_q(\mathfrak g)$ 在 $q=e^{2\pi\ii/(k+h^\vee)}$
处的表示范畴密切相关；这是 Kazhdan--Lusztig 等人工作的主题。

\textbf{第五，一般群。} 原文主要处理 $SU(N)$ 与 $SO(N)$；一般单群的情形类似，
但归一化与可积性条件需要逐群检查。

\subsection{两个常见的误读}

\subsubsection{CS/WZW 不是全息对偶}

第一个误读是把 CS/WZW 的对应说成“三维引力与二维 CFT 的全息对偶”，并由此联想到
AdS$_3$/CFT$_2$。两者不是一回事。

CS/WZW 是\textbf{态空间的等同}：三维理论在曲面 $\Sigma$ 上的 Hilbert 空间，
与二维理论在 $\Sigma$ 上的共形块空间，是同一个向量空间。这里没有“体”与“边界”的
动力学关系，没有引力子，没有全息字典。三维 Chern--Simons 理论的作用量 (2.1) 不含度规，
它根本没有引力自由度。

AdS$_3$/CFT$_2$ 是\textbf{动力学等价}：三维反德西特空间里的量子引力，与它边界上的
二维共形场论，描述同一套物理。它有引力子、有 BTZ 黑洞、有边界条件与全息字典。
两者唯一的接触点是：三维引力在适当条件下可以写成 Chern--Simons 理论
（规范群取非紧的 $SL(2,\RR)\times SL(2,\RR)$ 或其推广）。但这是一个“改写”，不是“等同”：
它要求非紧群、要求特定的边界条件，且只覆盖引力的经典/相空间层面。

读者若在文献里看到“CS 是全息的”，要追问：是哪一个 CS、哪一个群、哪一层级的对应。

\subsubsection{Khovanov 同调不是本文的预言}

第二个误读是把 Jones 多项式的范畴化（Khovanov 同调）说成“Witten 理论的自然推论”。
也不是。

Khovanov 同调是把 Jones 多项式从一个数提升为一个复形（或一个分次向量空间），
使得取 Euler 特征（带符号的维数）时退回 Jones 多项式。它满足函子性：链环的协边
给出复形之间的映射。这些性质说明它是一个“范畴化”的对象，比多项式携带更多信息。

但 1989 年的文章里没有任何关于同调、链复形或范畴化的暗示。把 Chern--Simons 理论
与 Khovanov 同调联系起来，要等到 Witten 2011 年（把 Khovanov 同调解释为
“高一维理论的边界条件”）以及 Kapustin--Witten 2007 年（4d 扭曲超对称规范理论的框架）。
即使在那之后，“CS 的 Wilson 线给出 Khovanov 复形”仍然不是一个已证明的定理；
它是一族猜想的统称。

\subsection{仍然开放的问题}

\begin{enumerate}[leftmargin=*]
  \item \textbf{函子性的严格化。} 把“缝合对应张量积”变成定理，需要什么条件？
        对一般的 $G$ 与一般的曲面，模函子的构造在什么范围内已知？
  \item \textbf{一般手术下的标架追踪。} (7.6) 在任意手术下的完整形式，
        以及它与缎带结构的关系。
  \item \textbf{半形式的系统处理。} $k\to k+h^\vee$ 在几何量子化框架里的一致推导，
        包括奇点处的相交上同调。
  \item \textbf{量子群与共形块的精确对应。} 可积表示范畴与 $U_q(\mathfrak g)$ 表示范畴
        在单位根处的等价，其精确陈述与证明。
  \item \textbf{与 Khovanov 同调的关系。} 是否存在一个“范畴化的 Chern--Simons 理论”，
        其边界给出 Khovanov 复形？
\end{enumerate}

\subsection{一条时间线}

把本文放进一条时间线里，能看清哪些是“当时的问题”、哪些是“后来的补课”。

\begin{center}
\small
\begin{tabular}{@{}p{1.9cm}p{7.4cm}p{4.6cm}@{}}
\toprule
年份 & 事件 & 与本文的关系 \\
\midrule
1984 & Jones 发现新的纽结不变量（von Neumann 代数） & 本文要解释的对象 \\
1985 & Jones 的 Bulletin AMS 短文；Kauffman 的状态和模型 & 组合定义成熟 \\
1986–87 & Drinfeld–Jimbo 量子群；Turaev–Reshetikhin 框架 & 提供二维/代数侧的语言 \\
1988 & Atiyah 在 Weyl 讨论会提两问；Witten 完成本文 & 本文的直接动机 \\
1989 & Moore–Seiberg、TUY 证明 Verlinde 公式（数学侧） & 补上本文只断言的部分 \\
1991 & Reshetikhin–Turaev 不变量；ADPW 的模空间量子化 & 严格化与几何化 \\
1994 & Faltings 证明 Verlinde 公式 & 数学证明的另一条路 \\
2000 & Khovanov 同调 & 范畴化，超出本文范围 \\
2007/2011 & Kapustin–Witten；Witten 的范畴化文章 & 范畴化的物理来源 \\
\bottomrule
\end{tabular}
\end{center}

\noindent 从这条线可以看出：本文处在“物理给出构造、数学随后补证明”的中间位置。
它不是一个封闭的定理，而是一个纲领；纲领的价值在于它指出了正确的对象
（曲面上的共形块），而不在于它把每一步都证完了。

\subsection{文献地图}

读到这里，读者可能需要知道去哪里查细节。下面按主题列出最常用的几份资料；
它们覆盖的范围、符号约定与严格程度各不相同，配合使用最有效。

\begin{center}
\small
\begin{tabular}{@{}p{3.0cm}p{5.7cm}p{5.4cm}@{}}
\toprule
资料 & 适合用来查什么 & 注意 \\
\midrule
Witten 1989（本文） & 原始论证、公式编号、物理动机 & 符号约定与教材不同；半形式略去 \\
河野《场论与拓扑学》 & 仿射代数、KZ、手术公式、Verlinde & CS 归一化差 $2\pi k$ \\
Schottenloher & Verlinde 的数学地位与证明史 & 第 11 章；物理推导不重复 \\
Moore TASI 2019 & 相空间量子化、标架反常、RCFT & 只查需要的节，不必通读 \\
Atiyah--Bott 1982 & 模空间的辛结构与 Kähler 结构 & 经典文献，符号较老 \\
TUY 1989 / Faltings 1994 & Verlinde 的数学证明 & 难度高，先读综述 \\
\bottomrule
\end{tabular}
\end{center}

\noindent 三条使用纪律：比较 $k$ 之前先对齐 $\Tr$ 与作用量的归一化；
看到 $q=e^{2\pi\ii/(k+h^\vee)}$ 时确认对方是否已做半形式平移；
凡涉及“标架”“相位”，先分清是配分函数的相位还是 Wilson 线的相位。

\subsection{开放问题的现状}

前面列的五个问题，到今天各有不同程度的进展，但没有一个被完全解决。
函子性方面，Segal 的框架与模函子理论给出了严格定义，但对一般 $G$ 与一般曲面，
构造仍依赖技术假设；手术公式的严格化（Reshetikhin--Turaev）覆盖了闭流形与链环，
但“缝合 = 张量积”的完整陈述仍不是一句话能证完的。
标架与半形式方面，两者的等价性在物理文献里是标准做法，
但一般奇点情形下的统一处理仍需要相交上同调与抛物丛的技术。
量子群与共形块的对应，在单位根处有 Kazhdan--Lusztig 等人的深刻结果，
但“可积表示范畴 $\cong$ $U_q(\mathfrak g)$ 的某子范畴”的精确陈述依赖 $q$ 的选取与范畴的完备化，
文献里的版本并不完全一致。范畴化方向是过去二十年最活跃的：
Khovanov 同调与 Khovanov--Rozansky 同调，以及它们与四维超对称规范理论的联系，
已经构成独立的研究领域；本讲义只讲到“它是什么、与本文什么关系”。

\subsection{怎么读这篇文章}

最后给三条读法，供不同背景的读者选择。

\textbf{快速读法（一天）。} 读前言、第一讲、第二讲、第七讲的小结。抓住四个词：
作用量、相空间、共形块、融合规则。

\textbf{标准读法（一周）。} 第一讲到第七讲按顺序读，每讲的公式自己推一遍。
第三、四讲是技术核心；第七讲是计算核心。

\textbf{研究读法（数周）。} 以原文为主，配合河野俊丈《场论与拓扑学》与 Schottenloher
《共形场论的数学导引》；遇到符号约定不一致时，回到原文页图核对。重点核对三件事：
$k$ 的归一化、$\Tr$ 的表示、半形式修正是否算进 $k$。

\subsection{结语}

回看全书，主线只有一条：一个不依赖度规的三维作用量，在曲面上量子化之后，
给出的态空间恰好是二维共形场论的共形块；纽结是这条主线上的标记点，
而 Jones 多项式是这条主线的可计算影子。

这条主线之所以值得写成一本书，是因为它同时是三个领域的交汇：
拓扑（三维流形的分类）、代数（量子群与融合范畴）、分析（几何量子化与模空间）。
任何一个领域单独看，都能看到一部分；只有把三者放在一起，才能理解为什么 Jones 多项式
会“出现”在三维理论里。

如果读者读完只记住一句话，我希望是这一句：\textbf{Chern--Simons 理论把“三维的拓扑”
编码成“二维的共形数据”，而量子化的曲面正是这两者之间的翻译词典。}

\begin{quote}
\small
\textbf{本讲要点。} 原文的五处空白（函子性、标架、半形式、量子群、一般群）是后续工作的入口；
CS/WZW 是态空间等同，不是全息对偶；Khovanov 同调是范畴化，不是本文的预言。
\end{quote}


\begin{thebibliography}{99}
\small

\bibitem{atiyah} M.~F. Atiyah,
\newblock New invariants of $3$- and $4$-manifolds,
\newblock {\it Proc.\ Symp.\ Pure Math.} {\bf 48} (1988), 285--299.

\bibitem{witten-tqft} E.~Witten,
\newblock Topological quantum field theory,
\newblock {\it Comm.\ Math.\ Phys.} {\bf 117} (1988), 353--386.

\bibitem{witten} E.~Witten,
\newblock Quantum field theory and the Jones polynomial,
\newblock {\it Comm.\ Math.\ Phys.} {\bf 121} (1989), 351--399.

\bibitem{jones} V.~F.~R. Jones,
\newblock A polynomial invariant for knots via von Neumann algebras,
\newblock {\it Bull.\ Amer.\ Math.\ Soc.} {\bf 12} (1985), 103--111;
\newblock Hecke algebra representations of braid groups and link polynomials,
\newblock {\it Ann.\ of Math.} {\bf 126} (1987), 335--388.

\bibitem{kauffman} L.~H. Kauffman,
\newblock State models and the Jones polynomial,
\newblock {\it Topology} {\bf 26} (1987), 395--407.

\bibitem{turaev} V.~G. Turaev,
\newblock The Yang--Baxter equation and invariants of links,
\newblock {\it Invent.\ Math.} {\bf 92} (1988), 527--553.

\bibitem{tuy} A.~Tsuchiya, K.~Ueno and Y.~Yamada,
\newblock Conformal field theory on universal family of stable curves with gauge symmetries,
\newblock {\it Adv.\ Stud.\ Pure Math.} {\bf 16} (1989), 297--372.

\bibitem{verlinde} E.~Verlinde,
\newblock Fusion rules and modular transformations in $2$d conformal field theory,
\newblock {\it Nucl.\ Phys.} {\bf B300} (1988), 360--376.

\bibitem{ms} G.~Moore and N.~Seiberg,
\newblock Polynomial equations for rational conformal field theories,
\newblock {\it Phys.\ Lett.} {\bf B212} (1988), 451--460;
\newblock Naturality in conformal field theory,
\newblock {\it Nucl.\ Phys.} {\bf B313} (1989), 16--40.

\bibitem{moore-tasi} G.~Moore,
\newblock TASI lectures on the Chern--Simons theory,
\newblock arXiv:1901.xxxxx（中译见本书附录所列资料）。

\bibitem{segaltqft} G.~Segal,
\newblock The definition of conformal field theory,
\newblock in {\it Topology, Geometry and Quantum Field Theory}, LMS Lecture Note Ser.\ 308,
Cambridge Univ.\ Press (2004), 421--577.

\bibitem{bpz} A.~A. Belavin, A.~M. Polyakov and A.~B. Zamolodchikov,
\newblock Infinite conformal symmetry in two-dimensional quantum field theory,
\newblock {\it Nucl.\ Phys.} {\bf B241} (1984), 333--380.

\bibitem{ab} M.~F. Atiyah and R.~Bott,
\newblock The Yang--Mills equations over Riemann surfaces,
\newblock {\it Phil.\ Trans.\ Roy.\ Soc.\ London} {\bf A308} (1982), 523--615.

\bibitem{quillen} D.~Quillen,
\newblock Determinants of Cauchy--Riemann operators over a Riemann surface,
\newblock {\it Funct.\ Anal.\ Appl.} {\bf 19} (1985), 31--34.

\bibitem{kac} V.~G. Kac,
\newblock {\it Infinite Dimensional Lie Algebras}, 3rd ed.,
Cambridge Univ.\ Press (1990).

\bibitem{hitchin} N.~J. Hitchin,
\newblock Flat connections and geometric quantization,
\newblock {\it Comm.\ Math.\ Phys.} {\bf 131} (1990), 347--380.

\bibitem{kz} V.~G. Knizhnik and A.~B. Zamolodchikov,
\newblock Current algebra and Wess--Zumino model in two dimensions,
\newblock {\it Nucl.\ Phys.} {\bf B247} (1984), 83--103.

\bibitem{rt} N.~Yu. Reshetikhin and V.~G. Turaev,
\newblock Invariants of $3$-manifolds via link polynomials and quantum groups,
\newblock {\it Invent.\ Math.} {\bf 103} (1991), 547--597.

\bibitem{faltings} G.~Faltings,
\newblock A proof for the Verlinde formula,
\newblock {\it J.\ Alg.\ Geom.} {\bf 3} (1994), 347--374.

\bibitem{khovanov} M.~Khovanov,
\newblock A categorification of the Jones polynomial,
\newblock {\it Duke Math.\ J.} {\bf 101} (2000), 359--426.

\bibitem{w11} E.~Witten,
\newblock Fivebranes and knots,
\newblock arXiv:1101.3216.

\bibitem{kw} A.~Kapustin and E.~Witten,
\newblock Electric-magnetic duality and the geometric Langlands program,
\newblock {\it Commun.\ Number Theory Phys.} {\bf 1} (2007), 1--236.

\bibitem{bh} J.~D. Brown and M.~Henneaux,
\newblock Central charges in the canonical realization of asymptotic symmetries,
\newblock {\it Comm.\ Math.\ Phys.} {\bf 104} (1986), 207--226.

\bibitem{maldacena} J.~Maldacena,
\newblock The large $N$ limit of superconformal field theories and supergravity,
\newblock {\it Adv.\ Theor.\ Math.\ Phys.} {\bf 2} (1998), 231--252.

\bibitem{witten-gravity} E.~Witten,
\newblock $2+1$ dimensional gravity as an exactly soluble system,
\newblock {\it Nucl.\ Phys.} {\bf B311} (1988), 46--78.

\bibitem{kono} 河野俊丈，《场论与拓扑学》，岩波书店 (1998)；中译注本 2026。

\bibitem{schottenloher} M.~Schottenloher,
\newblock {\it A Mathematical Introduction to Conformal Field Theory}, 2nd ed.,
Springer (2008)；中译注本 2026。

\bibitem{si} 李思，《量子力学与几何》讲义，第 5 章（几何量子化）；中译注本 2026。

\end{thebibliography}

\appendix
\section{符号与约定}

\begin{center}
\small
\begin{tabular}{@{}ll@{}}
\toprule
符号 & 含义 \\
\midrule
$M$ & 紧致、定向、无边三维流形 \\
$\Sigma$ & 闭定向曲面（$M$ 沿它切开） \\
$G$ & 紧致单 Lie 群，$d=\dim G$，$h^\vee$ 对偶 Coxeter 数（$SU(N)$ 时 $h^\vee=N$） \\
$A$ & $G$-联络，取值于 $\mathfrak g$ 的一形式 \\
$F$ & 曲率，$F=\dd A+A\wedge A$ \\
$k$ & Chern--Simons 耦合常数（整数） \\
$L$ & Chern--Simons 作用量 (2.1) \\
$W_R(C)$ & 沿 $C$、表示 $R$ 的 Wilson 线 \\
$\cM$ & $\Sigma$ 上平坦联络模空间 \\
$\cM_J$ & 复结构 $J$ 下的 Kähler 模空间 \\
$\cL$ & 行列式线丛（预量子化线丛） \\
$\cH_\Sigma$ & 沿 $\Sigma$ 切开的量子 Hilbert 空间 \\
$H^0(\cM_J,\cL^{\otimes k})$ & Kähler 极化下的量子态空间 \\
$R_\lambda$, $\lambda$ & 最高权为 $\lambda$ 的不可约表示 \\
$S_{ij}$ & 模变换矩阵（Verlinde 公式中的 $S$ 矩阵） \\
$N_{ijk}$ & 融合系数（三点函数维数） \\
\bottomrule
\end{tabular}
\end{center}

\noindent\textbf{三个容易混的约定。}
\begin{enumerate}[leftmargin=*]
  \item 本讲的 $L$ 与河野俊丈书里的 $\mathrm{CS}$ 差一个常数因子：
        $L=2\pi k\,\mathrm{CS}_{\text{河野}}$。比较 $k$ 之前先对齐归一化。
  \item 半形式修正把 $k$ 平移为 $k+h^\vee$。本讲凡写 $q=e^{2\pi\ii/(k+h^\vee)}$ 处
        都默认已含此修正；若只用到 $\cL^{\otimes k}$，则 $k$ 是“裸”的。
  \item 边注 \Wp{370} 指 Witten 论文的 CMP 印刷页；中译本的 PDF 页是重排的，
        不能直接换算，请用页边注“原文 p.N”定位。
\end{enumerate}

\section{练习提示}

\begin{enumerate}[leftmargin=*]
  \item \textbf{三叶结。} 对三叶结的标准投影连续用三次拆接关系，最后归到平凡纽结；
        注意每次光滑化后出现的链环分量数。
  \item \textbf{平凡链环。} 对 $n$ 个分量的平凡链环反复使用拆接关系，得到递推式
        $J_n=(-t^{1/2}-t^{-1/2})J_{n-1}$。
  \item \textbf{大规范变换。} 把 $g$ 写成 $S^3$ 上的映射，用度数公式
        \[
        \deg g=\frac{1}{24\pi^2}\int\Tr\left((g^{-1}\dd g)^3\right).
        \]
  \item \textbf{泊松括号。} 把 (3.1) 写成 $\int p^a\dot q^a$ 的形式，
        其中 $p,q$ 是 $A$ 的两个分量；正则括号即得 (3.2)。
  \item \textbf{截面计数。} 用 Riemann--Roch：$\dim H^0=\deg L+1-g$（$g=1$ 时 $\deg L=k$）；
        半形式把 $\cL^{\otimes k}$ 换成 $\cL^{\otimes k}\otimes K^{1/2}$，度数加 $h^\vee$ 的相应倍数。
  \item \textbf{$S$ 矩阵。} 对 $SU(2)$、$k=1$：$S=\frac{1}{\sqrt2}\begin{pmatrix}1&1\\1&-1\end{pmatrix}$，
        $S^2=1$ 直接验证。
  \item \textbf{$Z(S^3)$。} 它等于
        \[
        S_{00}=\sqrt{\frac{2}{k+2}}\,\sin\frac{\pi}{k+2}.
        \]
        取 $k=1,2,3$ 时分别为 $1/\sqrt2$、$1/2$、$\sqrt{2/5}\sin(\pi/5)$。
  \item \textbf{融合系数。} 用 (7.4) 取 $i=j=k=1/2$，$m=0$，再用 $S$ 的显式值。
\end{enumerate}

\section{原文—讲义对照表}

下表把 Witten 原文的各节与本书的讲次对应起来，方便读者两边对照。
页码是 CMP 印刷页；“公式”列给出该节的关键公式编号。

\begin{center}
\small
\begin{tabular}{@{}p{3.4cm}p{1.6cm}p{2.4cm}p{4.4cm}@{}}
\toprule
原文 & 页 & 本书 & 关键内容 \\
\midrule
引言 & 351–353 & 第 1 讲 & Atiyah 两问；Jones 缺三维定义 \\
\S1 作用量 & 353–356 & 第 2 讲 & (1.3) 作用量；(1.4) $k\in\ZZ$；(1.5)(1.6) \\
\S2 弱耦合 & 356–365 & 第 3、6 讲 & (2.16)(2.17) $\eta$ 修正；$q$ \\
\S2.1 纽结 & 362–365 & 第 5 讲 & 穿刺点与 Wilson 线 \\
\S3 正则量子化 & 366–374 & 第 3、4 讲 & (3.1)–(3.3) 辛形式与约束 \\
\S3.1 全纯观点 & 368–369 & 第 4 讲 & 极化与全纯截面 \\
\S3.2 平坦丛 & 369–371 & 第 4 讲 & 射影平坦联络 \\
\S3.3 Wilson 线 & 371–373 & 第 5 讲 & (3.4) 算符值约束；BWB \\
\S3.4 带标记球面 & 373–374 & 第 6、7 讲 & (3.5)(3.6)；融合系数 \\
\S4 可计算性 & 374–394 & 第 7 讲 & (4.28) 手术；(4.51) Verlinde \\
\S4.5 具体手术 & 388–393 & 第 7 讲 & (4.41) 标架；(4.51) \\
\S4.6 辫群 & 394 & 第 7 讲 & Jones 由迹给出 \\
\S5 物理应用 & 394–397 & 第 6、8 讲 & 共形块配对；两处误读 \\
\bottomrule
\end{tabular}
\end{center}

\section{术语表}

\begin{itemize}[leftmargin=*]
  \item \textbf{预量子化（prequantization）}：由整性条件 $[\omega]/2\pi\in H^2(M,\ZZ)$ 决定的线丛与联络。
  \item \textbf{极化（polarization）}：复化切丛的一个可积拉格朗日子丛；用来把 $L^2$ 截面砍成全纯截面。
  \item \textbf{半形式修正（metaplectic correction）}：额外乘 $K^{1/2}$，使能级平移 $k\to k+h^\vee$。
  \item \textbf{共形块（conformal block）}：带标记点曲面上 WZW 模型的态空间，也是 CS 理论在曲面上的态空间。
  \item \textbf{融合规则（fusion rule）}：两个表示张量积分解成不可约表示的重数 $N_{ijk}$。
  \item \textbf{$S$ 矩阵}：模群生成元 $S$ 在环面共形块上的作用，其矩阵元由 (7.2) 给出。
  \item \textbf{手术（surgery）}：挖去管状邻域再以不同方式粘回，用于构造三维流形。
  \item \textbf{标架反常（framing anomaly）}：配分函数在标架改变下的相位，由 $c$ 与标架整数决定。
  \item \textbf{范畴化（categorification）}：把多项式不变量提升为链复形/同调，Euler 特征退回原不变量。
\end{itemize}

\section{部分练习解答}

下面给出几道题的完整解答，其余留给读者；每道题的“关键一步”都写出来了，
卡住时对照 \S B 的提示。

\subsection*{第 1 讲第 1 题：三叶结}

三叶结的标准投影有三个正交叉。对其中一个交叉用 (1.1)：
\[
t^{-1}J(L_+)=t\,J(L_-)+(t^{1/2}-t^{-1/2})J(L_0).
\]
$L_-$ 是两个分量的平凡链环，其 Jones 多项式为 $-(t^{1/2}+t^{-1/2})$；
$L_0$ 是平凡纽结，值为 $1$。代回并整理，再对剩余交叉重复，得到
\[
J(\text{三叶结})=-t^{-4}+t^{-3}+t^{-1}.
\]
关键在于每一步都要重新判断 $L_-,L_0$ 的类型与分量数，这是最容易出错的地方。

\subsection*{第 3 讲第 4 题：泊松括号}

把 (3.1) 写成 $\int\dd t\int_\Sigma \epsilon^{ij}\Tr(A_i\partial_tA_j)$。
令两个空间分量为 $q^a=A_1^a$、$p^a=(k/4\pi)\epsilon^{ab}A_b$（这里 $a,b$ 是二维指标），
则作用量是 $\int p^a\dot q^a$ 的形式，正则括号 $\{q^a,p^b\}=\delta^{ab}$ 立即给出 (3.2)。
唯一需要注意的是 $\Tr$ 的归一化与 $\epsilon$ 的约定，两者都会影响系数 $4\pi/k$。

\subsection*{第 5 讲第 3 题：截面维数}

$G=SU(2)$、$\Sigma=T^2$、$M_J$ 为枕头套。取 $g=1$，Riemann--Roch 给
$\dim H^0(\cM,\cL^{\otimes k})=\deg(\cL^{\otimes k})=k$。
加上半形式 $K^{1/2}=\cL^{-2}$，度数减去 $2$ 的相应倍数后得 $k+1$。
这与“有效 level $k+2$ 的可积表示数减一”一致：
$k+2$ 个表示 $j=0,\frac12,\dots,\frac{k+1}{2}$，减去真空给出的那个平凡维数。

\subsection*{第 7 讲第 2 题：$Z(S^3)$}

$S^3$ 可以看作沿平凡纽结做手术（标架 $0$）得到，手术公式给 $Z(S^3)=S_{00}$。
代入 (7.2)：
\[
S_{00}=\sqrt{\frac{2}{k+2}}\sin\frac{\pi(2\cdot0+1)(2\cdot0+1)}{k+2}
=\sqrt{\frac{2}{k+2}}\sin\frac{\pi}{k+2}.
\]
$k=1$：$\sqrt{2/3}\sin(\pi/3)=1/\sqrt2$；$k=2$：$\sqrt{2/4}\sin(\pi/4)=1/2$。
这两个数可以用 (7.4) 与融合系数独立核对。

\subsection*{第 7 讲第 3 题：$k=1$ 的融合系数}

$k=1$ 时可积表示只有 $j=0,\frac12$。由 (7.4) 与 $S$ 的显式值，
$N_{1/2,1/2,0}=1$（$1/2\otimes1/2$ 含真空），$N_{1/2,1/2,1/2}=0$（$j=1$ 不可积）。
用 $S$ 的公式验证：
\[
N_{1/2,1/2,0}=\sum_{l=0,1/2}\frac{S_{1/2,l}S_{1/2,l}S_{0,l}}{S_{0,l}}
=\sum_l \frac{S_{1/2,l}^2}{1}=1,
\]
因为 $S_{1/2,0}^2=S_{1/2,1/2}^2=\frac12$。这正是 Verlinde 公式在最小情形的手算。

\section{复现与工具}

本讲义的源文件是 \texttt{preamble.tex}、\texttt{main.tex}、\texttt{chapters/L0--L8}、
\texttt{chapters/refs.tex} 与 \texttt{chapters/appendix.tex}。编译需要 xelatex
（ctex、fandol、tcolorbox、longtable、enumitem、booktabs、hyperref）：

\begin{quote}
\small\ttfamily
xelatex -interaction=nonstopmode main.tex\\
xelatex -interaction=nonstopmode main.tex
\end{quote}

\noindent 两遍是必要的：第一遍写目录与交叉引用，第二遍才排对。

同一仓库还包含另外三部分：
\begin{itemize}[leftmargin=*]
  \item 注疏版（按原文节序逐句对照）；
  \item 数据校验与锚点检查脚本；
  \item 协作生产流程记录（角色、议程、逐字稿、评分）。
\end{itemize}
\texttt{deliverables/workflow/} 记录本讲义的协作生产流程。

引用原文时，本讲义用 CMP 印刷页（351--399），可在 \S C 的对照表里查到每节位置。
中文译本重排为 40 页，不能直接按 PDF 页换算，请用页边注“原文 p.N”定位。


\end{document}
